% !TeX spellcheck = pt_BR
\documentclass[12pt,portuguese]{article}

%----------------------------------------------------------------------
% PACOTES ESSENCIAIS E TRADUÇÃO
%----------------------------------------------------------------------
\usepackage[brazilian]{babel}     % Suporte à língua portuguesa
\usepackage[utf8]{inputenc}       % Codificação de entrada
\usepackage[T1]{fontenc}          % Codificação da fonte
\usepackage{lmodern}              % Fonte escalável

%----------------------------------------------------------------------
% PACOTES MATEMÁTICOS
%----------------------------------------------------------------------
\usepackage{amsmath, amsthm, amssymb, amstext, mathtools}
\usepackage{mathrsfs}             % Para fontes como \mathscr
\usepackage{dsfont}               % Para o indicador de conjunto \mathds{1}
\usepackage[all]{xy}              % Para diagramas comutativos
\usepackage{imakeidx}             % Criação de Lista de Símbolos / Índice

% Desigualdades sempre com símbolos inclinados
\let\leq\leqslant
\let\le\leqslant
\let\ge\geqslant
\let\geq\geqslant

%----------------------------------------------------------------------
% FORMATAÇÃO, LAYOUT E CORES
%----------------------------------------------------------------------
\usepackage{geometry}
\usepackage{varwidth}
\usepackage{graphicx}
\usepackage[pdftex]{xcolor}
\usepackage{titlesec}
\usepackage{appendix}
\usepackage{enumitem}
\usepackage{float}
\usepackage{booktabs}
\usepackage{array}
\usepackage{colortbl}
\usepackage{fontawesome5}
\usepackage{lucide-icons}
\usepackage{caption}
\usepackage{stackengine}
\captionsetup{width=0.8\textwidth}

% Definição de Cores
\definecolor{teal}{RGB}{15,118,110}
\definecolor{text}{RGB}{28,25,23}
\definecolor{muted}{RGB}{87,83,78}
\definecolor{card}{RGB}{255,255,255}
\definecolor{border}{RGB}{231,229,228}

% Definição de margens e espaçamento
\setlength{\topmargin}{-0.7in}
\setlength{\oddsidemargin}{0in}
\setlength{\textheight}{9.55in}
\setlength{\textwidth}{6.5in}
\setlength{\baselineskip}{12mm}

% Formatação dos títulos de seção
\titleformat{\section}{\Large\bfseries}{\thesection.}{0.5em}{}
\titleformat{\subsection}{\large\bfseries}{\thesubsection.}{0.5em}{}
\titleformat{\subsubsection}{\normalfont\itshape}{\thesubsubsection.}{0.6em}{}

%----------------------------------------------------------------------
% TIKZ E GRÁFICOS
%----------------------------------------------------------------------
\usepackage{tikz}
\usepackage{pgfplots}
\pgfplotsset{compat=1.17}
\usetikzlibrary{positioning,calc,fit,backgrounds,nfold,decorations.pathreplacing,calligraphy,tikzmark,patterns,arrows.meta,bending,hobby}

%----------------------------------------------------------------------
% REFERÊNCIAS, LINKS E METADADOS DO PDF
%----------------------------------------------------------------------
\usepackage[pdftex, colorlinks, linkcolor={red!50!black}, 
citecolor={red}, urlcolor={blue!80!black}]{hyperref}

\hypersetup{
pdftitle={Introdução aos Processos Estocásticos e o Processo de Ramificação},
pdfauthor={L. Cioletti},
pdfsubject={processos estocásticos, processo de ramificação, funções
geradoras de probabilidade, soma aleatória, probabilidade de extinção},
pdfkeywords={processo de ramificação, processos estocásticos, função
geradora de probabilidade, soma aleatória, probabilidade de extinção,
equação de extinção, MAT-UnB} }

% --- PACOTE E COMANDO PARA A LICENÇA CREATIVE COMMONS ---
\usepackage[type={CC},modifier={by-nc-sa},version={4.0}]{doclicense}
\newcommand\blfootnote[1]{%
\begingroup
\renewcommand\thefootnote{}\footnote{#1}%
\addtocounter{footnote}{-1}%
\endgroup
}

%----------------------------------------------------------------------
% ALIASCNT E AMBIENTES DE TEOREMA (Cleveref)
%----------------------------------------------------------------------
\usepackage{aliascnt}
\theoremstyle{definition}
\newtheorem{theorem}{Teorema}

\newaliascnt{definition}{theorem}
\newtheorem{definition}[definition]{Definição}
\aliascntresetthe{definition}

\newaliascnt{example}{theorem}
\newtheorem{example}[example]{Exemplo}
\aliascntresetthe{example}

\newaliascnt{proposition}{theorem}
\newtheorem{proposition}[proposition]{Proposição}
\aliascntresetthe{proposition}

\newaliascnt{corollary}{theorem}
\newtheorem{corollary}[corollary]{Corolário}
\aliascntresetthe{corollary}

\newaliascnt{lemma}{theorem}
\newtheorem{lemma}[lemma]{Lema}
\aliascntresetthe{lemma}

\newaliascnt{remark}{theorem}
\newtheorem{remark}[remark]{Observação}
\aliascntresetthe{remark}

% O primeiro argumento é o espaço acima do ambiente. Acrescentamos 0,4 cm ao
% valor usual de 3pt para separar melhor os exercícios.
\newtheoremstyle{plainexercise}
  {14.4pt}{3pt}{\normalfont}{}{\bfseries}{.}{.5em}
  {\thmnumber{#2}\thmnote{ #3}}
\theoremstyle{plainexercise}
\newtheorem{exercise}{Exercício}

% Versão marcada dos exercícios: a estrela fica pendurada na margem esquerda,
% como nas notas anteriores.
\newtheoremstyle{starredexercise}
  {14.4pt}{3pt}{\normalfont}{}{\bfseries}{.}{.5em}
  {\llap{\textasteriskcentered\,}\thmnumber{#2}\thmnote{ #3}}
\theoremstyle{starredexercise}
\newtheorem{starredexercise}[exercise]{Exercício}
\theoremstyle{definition}

% Ambientes não numerados usados no texto
\newtheorem*{property}{Propriedade}
\newtheorem*{solution}{Solução}

% Se um enumerate for o primeiro conteúdo de um teorema/exercício (cabeçalho
% ainda pendente, \@inlabel true), o \leavevmode dispara o \everypar do amsthm
% que coloca o cabeçalho na própria linha, antes da lista iniciar.
\makeatletter
\setlist[enumerate]{before={\if@inlabel\leavevmode\fi}}
\makeatother

% Nomes em português para os ambientes
\AtBeginDocument{\renewcommand{\proofname}{\normalfont\textbf{Prova}}}
\renewcommand{\qedsymbol}{$\blacksquare$}

% Cleveref (deve vir depois de hyperref e amsthm)
\usepackage[brazilian,nameinlink]{cleveref}

\crefname{theorem}{teorema}{teoremas}
\Crefname{theorem}{Teorema}{Teoremas}
\crefname{definition}{definição}{definições}
\Crefname{definition}{Definição}{Definições}
\crefname{lemma}{lema}{lemas}
\Crefname{lemma}{Lema}{Lemas}
\crefname{proposition}{proposição}{proposições}
\Crefname{proposition}{Proposição}{Proposições}
\crefname{corollary}{corolário}{corolários}
\Crefname{corollary}{Corolário}{Corolários}
\crefname{example}{exemplo}{exemplos}
\Crefname{example}{Exemplo}{Exemplos}
\crefname{exercise}{exercício}{exercícios}
\Crefname{exercise}{Exercício}{Exercícios}
\crefname{remark}{observação}{observações}
\Crefname{remark}{Observação}{Observações}
\crefname{section}{seção}{seções}
\Crefname{section}{Seção}{Seções}

%----------------------------------------------------------------------
% MACROS ADICIONAIS
%----------------------------------------------------------------------
\let\oldemptyset\emptyset
\let\emptyset\varnothing
\graphicspath{{./}{../arquivos-aux/}}

% Seno na convenção em português
\renewcommand{\sin}{\operatorname{sen}}

%----------------------------------------------------------------------
% DOCUMENTO PRINCIPAL
%----------------------------------------------------------------------
\begin{document}

% --- CABEÇALHO COM O LOGO DA UnB ---
\vspace*{-0.8cm}
\begin{figure}[H]
\includegraphics{logo-mat-unb.pdf}
\end{figure}

% --- TÍTULO, AUTOR E DATA ---
\vspace*{0.3cm}
\begin{center}
{\LARGE\bf Introdução ao Processo de Ramificação}
\vspace{1cm}

{\large \bf L. Cioletti}
\vspace{0.3cm}

{\footnotesize Novembro de 2026}
\end{center}

% --- RESUMO (ABSTRACT) ---
\begin{abstract}
Nestas notas, iniciamos o estudo de processos que evoluem aleatoriamente ao
longo do tempo e desenvolvemos o primeiro exemplo completo dessa classe: o
processo de ramificação, modelo simples para o crescimento de uma população
que se reproduz. Apresentamos o modelo com tamanhos de família independentes
e identicamente distribuídos e mostramos como a função geradora de
probabilidade do tamanho da $n$-ésima geração é obtida como a $n$-ésima iterada da função
geradora do tamanho de família. Obtemos também o valor médio da população neste modelo
e analisamos seu comportamento assintótico. Resolvemos em forma fechada o
caso em que os tamanhos de família têm distribuição geométrica, determinamos
a probabilidade de extinção como a menor raiz não negativa da equação
$x=G(x)$ e demonstramos que a extinção é certa se e somente se
$\mu\leqslant 1$. Para tornar a exposição autossuficiente, reapresentamos com
provas os instrumentos necessários sobre funções geradoras de probabilidade,
em particular a fórmula da soma aleatória. O texto conta com duas figuras, a
árvore genealógica do processo e as interseções de $y=G(x)$ com $y=x$, e com
exercícios e problemas propostos.
A exposição segue de perto as Seções 9.1 a 9.6 do Capítulo 9 da
referência \cite{GrimmettWelsh2014}.
\end{abstract}

% --- LICENÇA NO RODAPÉ ---
\blfootnote{%
\doclicenseImage[imagewidth=2.6em]\
\begingroup
\hypersetup{urlcolor=darkgray}
\href{\doclicenseURL}{%
\doclicenseName{} -- Attribution-NonCommercial-ShareAlike 4.0 International%
}%
\endgroup
}

% Origem: G&W, seção 9.1 (Random processes).
\section{Introdução}\label{sec:introducao}

Nas notas de aula anteriores desenvolvemos diversos dos resultados
básicos da Teoria de probabilidade.
Nestas notas iniciamos o estudo de aplicações dessa teoria
a fenômenos que evoluem aleatoriamente ao longo do tempo.

A passagem do tempo faz parte do mundo em que vivemos
e, consequentemente,
muitas aplicações da probabilidade envolvem quantidades
que se desenvolvem de modo aleatório à medida que o tempo passa.
Processos que evoluem aleatoriamente dessa forma são chamados de
\emph{processos aleatórios} ou \emph{processos estocásticos},
e há muitos tipos diferentes deles.
A maior parte dos processos reais da natureza,
como a contagem de pólen em uma cidade
ou a posição de um time no campeonato brasileiro de futebol,
evolui segundo regras complicadas demais
para serem descritas com exatidão,
e bons modelos probabilísticos para esses processos
podem ser bastante complexos.
Até o final do curso vamos nos restringir a alguns dos processos aleatórios mais simples,
entre eles vamos considerar com mais profundidade 
\begin{enumerate}[label=(\alph*)]
  \item \emph{processos de ramificação}: modelam o crescimento de uma
        população que se reproduz por si mesma (como a população humana);
  \item \emph{passeios aleatórios}: modelam o movimento de uma partícula
        que se desloca de modo errático dentro de um meio (uma partícula
        de poeira na atmosfera, por exemplo);
  \item \emph{processos de Poisson e processos relacionados}: modelam
        fenômenos como a emissão de partículas radioativas por uma fonte
        que decai lentamente, ou o comprimento da fila no caixa de um
        supermercado.
\end{enumerate}
Há uma teoria razoavelmente completa de cada um desses três tipos de
processo. Nestas notas trataremos de alguns aspectos mais elementares
da teoria do primeiro tipo, o processo de ramificação, 
que pode ser usado como um modelo para o crescimento de populações. 

% Origem: reapresentação condensada das notas func-gen-discretas (G&W, cap. 4).
\section{Preliminares sobre Funções Geradoras}\label{sec:preliminares-fgp}

O processo de ramificação que vamos estudar evolui por uma regra simples: cada
indivíduo de uma geração deixa um número aleatório de filhos, e o tamanho da
geração seguinte é a soma desses números. O objeto central destas notas é,
portanto, uma soma com um número aleatório de parcelas, e a ferramenta natural
para analisá-la é a função geradora de probabilidade. 

Para conveniência do leitor revisamos alguns dos conceitos e resultados fundamentais 
sobre funções geradoras necessários para o desenvolvimento do texto. 
Para maiores detalhes o leitor pode consultar as notas
de aula sobre \href{https://mat.unb.br/~cioletti/ensino/func-gen-discretas.pdf}{\emph{Funções
Geradoras de Probabilidade}}.

\begin{definition}\label{def:fgp}
Seja $X$ uma variável aleatória tomando valores em $\{0,1,2,\ldots\}$ e para cada
$k\in\{0,1,2,\ldots\}$ seja $p_k\equiv p_X(k)\equiv \mathbb{P}(X=k)$. A \emph{função
geradora de probabilidade} de $X$ é a função
definida por
\begin{equation}\label{eq:def-fgp}
G_X(s)\equiv \mathbb{E}[s^X] =\sum_{k=0}^{\infty}s^k p_k,
\end{equation}
para cada valor de $s\in\mathbb{R}$, onde a série acima converge absolutamente. 
\end{definition}

A normalização da função de probabilidade garante um domínio comum a todas as
funções geradoras de probabilidade. 
Mais precisamente, para qualquer variável aleatória
tomando valores em $\{0,1,2,\ldots\}$ temos
\begin{equation}\label{eq:fgp-convergencia}
\sum_{k=0}^{\infty}|p_k s^k|
=\sum_{k=0}^{\infty}p_k|s|^k
\leqslant\sum_{k=0}^{\infty}p_k=1
\qquad \forall s\in [-1,1]. 
\end{equation}
Logo, toda função geradora de probabilidade
converge absolutamente para $|s|\leqslant 1$ e, vista como série
de potências, tem raio de convergência pelo menos um. Em particular,
$G_X$ é contínua em $[0,1)$, pois toda série de potências é contínua no
interior do seu disco de convergência. Esse argumento, porém, não alcança
o ponto $s=1$: se o raio de convergência for exatamente um, esse ponto fica
na fronteira do disco, e a continuidade de uma série de potências não se
estende automaticamente à fronteira. Para examinar o comportamento de $G_X$
em $s=1$, enunciamos e provamos a seguir um resultado clássico sobre séries
com coeficientes não negativos, o lema de Abel.

% Origem: G&W, p. 55 (lema de Abel).
\begin{lemma}[Lema de Abel]\label{lem:abel}
Seja $(a_k)_{k\geqslant 0}$ uma sequência de números reais não negativos
tal que a série de potências $\sum_{k=0}^{\infty}a_ks^k$ convirja para
cada $s$ com $|s|<1$, e denote por $A(s)$ a sua soma. Então
\[
\lim_{s\uparrow 1}A(s)=\sum_{k=0}^{\infty}a_k
\quad\text{em }[0,+\infty],
\]
isto é, os dois lados são iguais, com a possibilidade de ambos valerem
$+\infty$.
\end{lemma}

\begin{proof}
A função $A$ é não decrescente em $[0,1)$: se $0\leqslant s\leqslant t<1$,
então $s^k\leqslant t^k$ para todo $k$ e, multiplicando por $a_k\geqslant 0$
e somando, obtemos
\[
\sum_{k=0}^{n}a_ks^k\leqslant\sum_{k=0}^{n}a_kt^k
\qquad\text{para todo } n,
\]
e, fazendo $n\to\infty$ e usando que as somas parciais convergem a $A(s)$ e
a $A(t)$, concluímos que $A(s)\leqslant A(t)$. Como toda função não decrescente admite
limite à esquerda como um único ponto em $[0,+\infty]$, podemos afirmar que 
está bem-definido $L=\lim_{s\uparrow 1}A(s)$.
Além do mais este limite coincide com o supremo dos valores $A(s)$, tomado sobre $0\leqslant s<1$.

Para ver que $\sum_{k=0}^{\infty}a_k\leqslant L$, fixe $n$. Como
$a_k\geqslant 0$, a soma parcial de ordem $n$ não ultrapassa a soma
total:
\[
\sum_{k=0}^{n}a_ks^k\leqslant A(s)
\qquad\text{para todo } s\in[0,1).
\]
O lado esquerdo, sendo um polinômio em $s$, converge a
$\sum_{k=0}^{n}a_k$ quando $s\uparrow 1$; fazendo $s\uparrow 1$ na
desigualdade, obtemos
\[
\sum_{k=0}^{n}a_k\leqslant L
\qquad\text{para todo } n.
\]
Como os $a_k$ são não negativos, as somas parciais $\sum_{k=0}^{n}a_k$
são não decrescentes em $n$. Além do mais supremo destas somas é dado por 
$\sum_{k=0}^{\infty}a_k$. Como a desigualdade acima, é valida para todo $n\in\mathbb{N}$
ela também é válida para o supremo e portanto
\[
\sum_{k=0}^{\infty}a_k\leqslant L.
\]

Para a desigualdade oposta, fixe $s\in[0,1)$. Já que $s^k\leqslant 1$ segue
\[
\sum_{k=0}^{n}a_ks^k\leqslant\sum_{k=0}^{n}a_k\leqslant\sum_{k=0}^{\infty}a_k
\qquad\text{para todo } n,
\]
e, fazendo $n\to\infty$, obtemos $A(s)\leqslant\sum_{k=0}^{\infty}a_k$.
Isso vale para todo $s\in[0,1)$ e, como $L$ é o supremo dos valores
$A(s)$, concluímos que $L\leqslant\sum_{k=0}^{\infty}a_k$, o que encerra a
demonstração.
\end{proof}

Observe que no caso de uma função geradora $G_X$, os coeficientes $p_k$ são não negativos
e a série converge absolutamente para $|s|\leqslant 1$, de modo que a
hipótese do \Cref{lem:abel} é satisfeita. Neste caso o lema garante que 
\[
\lim_{s\uparrow 1}G_X(s)=\sum_{k=0}^{\infty}p_k=1=G_X(1),
\]
e concluímos que $G_X$ é contínua à esquerda em $s=1$ e, portanto,
contínua em todo o intervalo $[0,1]$.

Dois valores de $s$ já revelam informação probabilística: em $s=0$, a série
se reduz ao coeficiente constante, que é a probabilidade de $X$ ser nulo;
em $s=1$, somamos todas as probabilidades. O ponto $s=0$ reaparecerá nas
seções seguintes: no processo de ramificação, uma geração vazia é uma
geração extinta, e as probabilidades de extinção da população serão
valores de funções geradoras nesse ponto.

\begin{proposition}\label{prop:fgp-basica}
Se $X$ tem função de probabilidade $(p_k)_{k\geqslant 0}$ e 
função geradora de probabilidade $G_X$, então
\[
G_X(0)=p_0=\mathbb{P}(X=0)
\quad\text{e}\quad
G_X(1)=1.
\]
\end{proposition}

\begin{proof}
As duas identidades seguem diretamente de \eqref{eq:def-fgp}: para $s=0$, os
termos com $k\geqslant 1$ se anulam; para $s=1$, obtemos
$\sum_{k=0}^{\infty}p_k=1$.
\end{proof}

O valor em $0$ extrai apenas $p_0$. Para recuperar os demais, derivamos a função geradora, 
já que séries de potências podem ser derivadas
termo a termo no interior do raio de convergência, de modo que, para todo $k\geqslant 0$,
temos
\[
\mathbb{P}(X=k)=p_k=\frac{G_X^{(k)}(0)}{k!},
\]
com a convenção $G_X^{(0)}=G_X$. Assim, podemos recuperar separadamente qualquer
entrada da função de probabilidade a partir da função geradora de probabilidade.

\medskip

O tamanho de cada geração do processo de ramificação é a soma dos números
de filhos deixados pelos indivíduos da geração anterior, de modo que
precisamos saber calcular a função geradora de uma soma de variáveis
aleatórias. Começamos com duas parcelas independentes, $X$ e $Y$, e
perguntamos: que operação sobre $G_X$ e $G_Y$ corresponde à soma $X+Y$?

O método direto usa a função de probabilidade: o evento $\{X+Y=k\}$ é a
união dos eventos $\{X=j\}\cap\{Y=k-j\}$, com $j=0,\ldots,k$, que são
dois a dois disjuntos; sob independência, as probabilidades conjuntas
fatoram, e a soma sobre $j$ produz a convolução
\[
\mathbb{P}(X+Y=k)=\sum_{j=0}^{k}\mathbb{P}(X=j)\,\mathbb{P}(Y=k-j)
\qquad\text{para } k=0,1,2,\ldots.
\]
O custo é uma soma nova para cada valor de $k$.

A definição da função geradora sugere um atalho: nela, cada probabilidade
é o coeficiente de uma potência de $s$, e o produto de potências de uma
mesma base soma os expoentes, $s^js^\ell=s^{j+\ell}$. Ao multiplicar um
termo genérico $\mathbb{P}(X=j)s^j$ de $G_X$ por um termo genérico
$\mathbb{P}(Y=\ell)s^\ell$ de $G_Y$, obtemos
$\mathbb{P}(X=j)\mathbb{P}(Y=\ell)s^{j+\ell}$; o coeficiente de $s^k$ no
produto $G_X(s)G_Y(s)$ reúne, portanto, os pares $(j,\ell)$ tais que
$j+\ell=k$, ponderados pelos mesmos pesos da convolução, e é exatamente
$\mathbb{P}(X+Y=k)$. Somar variáveis aleatórias independentes é
multiplicar as suas funções geradoras.

A prova formal é ainda mais curta: a mesma identidade, com as variáveis
aleatórias no lugar dos expoentes, é $s^{X+Y}=s^Xs^Y$, e a independência
fatora a esperança do produto.

\begin{theorem}[Produto de funções geradoras]\label{thm:fgp-produto}
Sejam $X$ e $Y$ variáveis aleatórias discretas independentes, ambas com
valores em $\{0,1,2,\ldots\}$. Então, para todo $s$ com $|s|\leqslant 1$,
\[
G_{X+Y}(s)=G_X(s)G_Y(s).
\]
\end{theorem}

\begin{proof}
Como $X$ e $Y$ são independentes, as variáveis $s^X$ e $s^Y$ também são
independentes. Para $|s|\leqslant 1$, ambas são limitadas em módulo por $1$
e, portanto, integráveis.
Logo,
\begin{align*}
G_{X+Y}(s)
&=\mathbb{E}[s^{X+Y}]
=\mathbb{E}[s^Xs^Y]
\\[0.2cm]
&=\mathbb{E}[s^X]\,\mathbb{E}[s^Y]
=G_X(s)G_Y(s),
\end{align*}
o que encerra a demonstração.
\end{proof}

A aplicação repetida do teorema revela a vantagem sobre a convolução.

\begin{corollary}\label{cor:fgp-soma-finita}
Sejam $X_1,\ldots,X_n$ variáveis aleatórias discretas independentes, com valores
em $\{0,1,2,\ldots\}$, e defina $S_n=X_1+\cdots+X_n$. Então
\[
G_{S_n}(s)=\prod_{i=1}^{n}G_{X_i}(s),
\qquad |s|\leqslant 1.
\]
Em particular, se as parcelas são identicamente distribuídas, com a mesma
distribuição de uma variável $X$, então
\begin{equation}\label{eq:fgp-soma-finita-iid}
G_{S_n}(s)=\bigl(G_X(s)\bigr)^n,
\qquad |s|\leqslant 1.
\end{equation}
\end{corollary}

\begin{proof}
O argumento é por indução em $n$; o caso $n=1$ é imediato. No passo
indutivo, a independência da família garante que
$S_{n-1}$ e $X_n$ são independentes, pois $S_{n-1}=X_1+\cdots+X_{n-1}$ é
função de $X_1,\ldots,X_{n-1}$, que são independentes de $X_n$. Aplicando
o \Cref{thm:fgp-produto} ao par $(S_{n-1},X_n)$ e usando a hipótese de
indução, obtemos
\[
G_{S_n}(s)
=G_{S_{n-1}}(s)G_{X_n}(s)
=\cdots
=\prod_{i=1}^{n}G_{X_i}(s).
\]
No caso identicamente distribuído, todos os fatores desse produto são
iguais a $G_X(s)$.
\end{proof}

\medskip

A função geradora de probabilidade também permite calcular momentos por derivação. Cada derivação de $s^k$
reduz o expoente em uma unidade, produzindo os fatores $k,k-1,\ldots$, e não a
potência $k^r$. Por isso, escrevemos $(X)_r=X(X-1)\cdots(X-r+1)$ para a
\emph{potência fatorial decrescente} de ordem $r$ e chamamos $\mathbb{E}[(X)_r]$
de \emph{$r$-ésimo momento fatorial} de $X$, admitindo o valor $+\infty$.
Derivar termo a termo uma série infinita e alcançar a fronteira $s=1$ exigem
justificativas distintas; o próximo resultado separa essas duas passagens.

\begin{theorem}\label{thm:fgp-momentos}
Seja $X$ uma variável aleatória com valores em $\{0,1,2,\ldots\}$ e seja $G_X$
sua função geradora de probabilidade. Para todo inteiro $r\geqslant 1$ e todo $|s|<1$,
\[
G_X^{(r)}(s)=\sum_{k=r}^{\infty}(k)_r\,s^{k-r}\,\mathbb{P}(X=k).
\]
Além disso, as derivadas laterais em $1$, definidas por
$G_X^{(r)}(1)\equiv\lim_{s\uparrow 1}G_X^{(r)}(s)$, satisfazem
$G_X^{(r)}(1)=\mathbb{E}[(X)_r]$, com a possibilidade de ambos os lados serem
iguais a $+\infty$. Em particular, $G'_X(1)=\mathbb{E}[X]$ e
$G''_X(1)=\mathbb{E}[X(X-1)]$, e, quando $\mathbb{E}[X^2]<\infty$,
\[
\operatorname{Var}(X)=G''_X(1)+G'_X(1)-\bigl(G'_X(1)\bigr)^2.
\]
\end{theorem}

\begin{proof}
A prova tem duas etapas. Primeiro justificamos a derivação no interior do
intervalo de convergência; depois passamos de $s<1$ para $s=1$.

Como $G_X(s)=\sum_{k=0}^{\infty}\mathbb{P}(X=k)s^k$ converge em $s=1$,
seu raio de convergência é pelo menos $1$. Pelo teorema de derivação
termo a termo de séries de potências, no interior do raio podemos derivar
termo a termo, e a série derivada tem o mesmo raio. Como $|s|<1$ está no
interior desse raio, podemos repetir o procedimento $r$ vezes. A $r$-ésima
derivada de $s^k$ é nula para $k<r$ e é $(k)_r s^{k-r}$ para
$k\geqslant r$, o que produz a fórmula enunciada.

Resta analisar o que acontece em $s=1$. Não podemos simplesmente substituir $s=1$ em uma
identidade demonstrada apenas para $|s|<1$. Os coeficientes
$(k)_r\,\mathbb{P}(X=k)$ da série derivada são não negativos e, como vimos
acima, essa série converge para cada $s$ com $|s|<1$, pois tem o mesmo raio
de convergência da original. O \Cref{lem:abel} aplica-se, portanto, e
fornece
\begin{align*}
\lim_{s\uparrow 1}G_X^{(r)}(s)
&=\sum_{k=r}^{\infty}(k)_r\,\mathbb{P}(X=k)
\\[0.3cm]
&=\mathbb{E}[(X)_r].
\end{align*}
A igualdade admite o valor $+\infty$ em ambos os lados, de modo que a
conclusão permanece válida mesmo quando o momento fatorial não é finito.

Os casos $r=1$ e $r=2$ fornecem as identidades anunciadas. Como
$(X)_1=X$, temos $\mathbb{E}[X]=G'_X(1)$; como $(X)_2=X(X-1)$, temos
$\mathbb{E}[X(X-1)]=G''_X(1)$. A derivada segunda produz, portanto, o
segundo momento fatorial, e não diretamente $\mathbb{E}[X^2]$. A
identidade $X^2=X(X-1)+X$ faz a conversão entre as duas formas:
\[
\mathbb{E}[X^2]=G''_X(1)+G'_X(1).
\]
Nas notas de aula sobre esperança discreta, obtivemos a expressão
$\operatorname{Var}(X)=\mathbb{E}[X^2]-(\mathbb{E}[X])^2$. Substituindo
nela as duas fórmulas anteriores, chegamos a
\[
\operatorname{Var}(X)
=G''_X(1)+G'_X(1)-\bigl(G'_X(1)\bigr)^2,
\]
o que encerra a demonstração.
\end{proof}

Essas identidades são fórmulas para momentos finitos. Se a derivada lateral em
$1$ vale $+\infty$, a afirmação sobre a esperança ou o momento fatorial continua
válida no sentido estendido, mas não devemos formar diferenças indeterminadas
para calcular a variância; seu uso pressupõe $G''_X(1)<\infty$.

\medskip

O produto de funções geradoras descreve uma soma com número fixo de parcelas;
no processo de ramificação, porém, o próprio número de parcelas é aleatório.
Nesse caso, fixamos esse número, usamos o produto e reunimos os casos possíveis,
o que transforma o produto em uma composição de funções geradoras. Adotamos a
convenção da soma vazia: $X_1+\cdots+X_N=0$ em $\{N=0\}$, com contribuição
$s^0=1$ para a função geradora de probabilidade, inclusive quando $s=0$
(convenção $0^0=1$).

\begin{theorem}[Fórmula da soma aleatória]\label{thm:soma-aleatoria}
Sejam $N,X_1,X_2,\ldots$ variáveis aleatórias com valores em
$\{0,1,2,\ldots\}$, com $X_1,X_2,\ldots$ independentes e identicamente
distribuídas, com a mesma distribuição de uma variável $X$ de função geradora
de probabilidade $G_X$, e $N$ independente da sequência $X_1,X_2,\ldots$. Defina
$S_N=X_1+\cdots+X_N$, com
$S_N=0$ em $\{N=0\}$. Então, para $|s|\leqslant 1$,
\begin{equation}\label{eq:fgp-soma-aleatoria}
G_{S_N}(s)=G_N\bigl(G_X(s)\bigr).
\end{equation}
\end{theorem}

\begin{proof}
Fixe $s$ com $|s|\leqslant 1$. Os eventos $\{N=n\}$, com
$n=0,1,2,\ldots$, são dois a dois disjuntos e sua união é o espaço
amostral. Em cada evento $\{N=n\}$, a soma aleatória se reduz à soma fixa
$X_1+\cdots+X_n$ (vazia, e igual a $0$, quando $n=0$). Portanto,
pontualmente,
\[
s^{S_N}
=\sum_{n=0}^{\infty}s^{X_1+\cdots+X_n}\,\mathds{1}_{\{N=n\}}.
\]
Para $0\leqslant s\leqslant 1$, todos os termos são não negativos, e a
convergência monótona permite trocar a esperança com a soma. Para
$-1\leqslant s<0$, a mesma troca é justificada por convergência dominada:
as somas parciais são limitadas em módulo por
$\sum_{n=0}^{\infty}\mathds{1}_{\{N=n\}}=1$. Em ambos os casos,
\begin{align*}
G_{S_N}(s)
&=\mathbb{E}\bigl[s^{S_N}\bigr]
\\[0.2cm]
&=\sum_{n=0}^{\infty}
\mathbb{E}\bigl[s^{X_1+\cdots+X_n}\,\mathds{1}_{\{N=n\}}\bigr].
\end{align*}
Essa é a decomposição sobre a partição $\{N=n\}$, escrita explicitamente
com indicadores; ela corresponde ao Teorema da Partição, apresentado nas
notas de aula sobre esperança condicional.

A independência entre $N$ e a sequência das parcelas garante que
$\mathds{1}_{\{N=n\}}$ é independente de $s^{X_1+\cdots+X_n}$, que é
função de $X_1,\ldots,X_n$. Para $|s|\leqslant 1$, ambas as
variáveis são limitadas em módulo por $1$, logo integráveis, e a esperança
do produto fatora:
\[
\mathbb{E}\bigl[s^{X_1+\cdots+X_n}\,\mathds{1}_{\{N=n\}}\bigr]
=\mathbb{E}\bigl[s^{X_1+\cdots+X_n}\bigr]\,\mathbb{P}(N=n).
\]

Em seguida, calculamos a função geradora de probabilidade da soma fixa. As
parcelas são independentes e, para $|s|\leqslant 1$, as variáveis $s^{X_i}$ são
limitadas
em módulo por $1$; o \Cref{cor:fgp-soma-finita}, enunciado nesse mesmo
domínio, aplica-se portanto diretamente. Como as parcelas têm a mesma
distribuição, todos os fatores são iguais:
\[
\mathbb{E}\bigl[s^{X_1+\cdots+X_n}\bigr]
=\prod_{i=1}^{n}\mathbb{E}\bigl[s^{X_i}\bigr]
=\bigl(G_X(s)\bigr)^n.
\]
Para $n=0$, a soma vazia e o produto vazio valem ambos $1$. Aqui aparecem
os papéis distintos das hipóteses: a independência entre as parcelas
produz o produto, enquanto a distribuição comum o transforma na potência
de uma única função.

Resta reunir os valores possíveis de $N$:
\begin{align}
G_{S_N}(s)
&=\sum_{n=0}^{\infty}\bigl(G_X(s)\bigr)^n\,\mathbb{P}(N=n)
\label{eq:fgp-soma-aleatoria-casos}
\\[0.2cm]
&=G_N\bigl(G_X(s)\bigr).\notag
\end{align}
A última igualdade é exatamente \eqref{eq:def-fgp}, a definição de $G_N$,
avaliada no ponto $G_X(s)$. Essa avaliação é legítima: para
$|s|\leqslant 1$, os coeficientes não negativos de $G_X$ dão
$|G_X(s)|\leqslant G_X(|s|)\leqslant G_X(1)=1$, de modo que $G_X(s)$
pertence ao domínio em que $G_N$ converge absolutamente. Assim, decompor sobre os
valores de $N$, aplicar o produto em cada caso e reunir os casos faz a
composição surgir, o que encerra a demonstração.
\end{proof}

É essa identidade que alimentará a recursão das funções geradoras nas seções
seguintes, com o número de parcelas dado pelo tamanho da geração anterior.

\section{Um Modelo para o Crescimento de Populações}\label{sec:modelo-populacao}
% Origem: G&W, seção 9.2 (A model for population growth).

Antes de usar a fórmula da soma aleatória, precisamos construir o modelo
populacional que ela descreve. No tempo $n=0$, existe apenas um indivíduo.
Cada indivíduo vive por uma unidade de tempo. Ao morrer, deixa zero ou mais
filhos. Os filhos nascidos no mesmo instante formam uma geração. A
\Cref{fig:ciclo-de-vida} ilustra como ocorre a evolução temporal do modelo.

\tikzset{
  vertice ramificacao/.style={circle, fill=text, inner sep=0pt, minimum size=2.2mm},
  aresta ramificacao/.style={draw=text, line width=0.6pt},
  eixo tempo/.style={draw=muted, line width=0.45pt},
  rotulo tempo/.style={font=\scriptsize, text=muted, inner sep=0pt},
  guia tempo/.style={draw=muted, opacity=0.19, line width=0.35pt, densely dashed},
  caixa familia/.style={draw=teal!65, fill=teal!9, line width=0.45pt,
                        rounded corners=1.2pt},
}

\newcommand{\arvoreciclodevida}[1]{%
  % Cada vértice é dado pelo nome, pelo pai, pela geração e pela abscissa.
  \foreach \nome/\pai/\geracao/\posicao in {
    g0/raiz/0/0,
    g1a/g0/1/-2, g1b/g0/1/2,
    g2a/g1a/2/-3.8, g2b/g1a/2/-2, g2c/g1a/2/-0.2, g2d/g1b/2/2,
    g3a/g2b/3/-2.8, g3b/g2b/3/-1.2, g3c/g2c/3/-0.2,
    g4a/g3a/4/-3.4, g4b/g3a/4/-2.2, g4c/g3b/4/-1.2, g4d/g3c/4/-0.2,
    g5a/g4b/5/-2.2, g5b/g4d/5/-0.75, g5c/g4d/5/0.35
  } {
    \ifnum\geracao>#1\relax
    \else
      \def\opacidade{0.16}
      \def\espessura{0.6pt}
      \ifnum\geracao=#1\relax
        \def\opacidade{1}
        \def\espessura{0.8pt}
      \fi
      \node[vertice ramificacao, opacity=\opacidade]
        (\nome) at (\posicao, -\geracao) {};
      \ifnum\geracao>0\relax
        \draw[aresta ramificacao, line width=\espessura, opacity=\opacidade]
          (\pai) -- (\nome);
      \fi
    \fi
  }%
}

\begin{figure}[H]
  \centering
  \begin{tikzpicture}[x=1cm, y=1.26cm, >=Stealth]
    \def\distanciacolunas{7.25}
    \def\eixodireito{2.5}
    \foreach \instante/\coluna/\linha in {0/0/0, 1/\distanciacolunas/0,
                                         2/0/-3.85, 3/\distanciacolunas/-3.85,
                                         4/0/-10.15, 5/\distanciacolunas/-10.15} {
      \begin{scope}[xshift=\coluna cm, yshift=\linha cm, name prefix=t\instante-]
        % Os eixos ficam nas bordas externas de cada par.
        \ifodd\instante\relax
          \let\ultimotempo\instante
          \let\eixo\eixodireito
          \def\pontadamarca{2.35}
          \def\ancora{west}
        \else
          \def\eixo{-4.3}
          \def\pontadamarca{-4.15}
          \def\ancora{east}

          % As guias de cada par são desenhadas antes de ambas as árvores.
          \pgfmathtruncatemacro{\ultimotempo}{\instante+1}
          \foreach \nivel in {0,...,\ultimotempo}
            \draw[guia tempo] (\eixo, -\nivel) --
              ({\distanciacolunas+\eixodireito}, -\nivel);
        \fi
        \draw[eixo tempo, ->] (\eixo, 0.3) --
          ([yshift=-2mm]\eixo, {-\ultimotempo-0.25});
        \foreach \nivel in {0,...,\ultimotempo} {
          \draw[eixo tempo] (\eixo, -\nivel) -- (\pontadamarca, -\nivel);
          \ifnum\nivel=\instante\relax
            \node[rotulo tempo, font=\small, text=text, anchor=\ancora, inner xsep=2.5pt]
              at (\eixo, -\nivel) {$n=\instante$};
          \else
            \node[rotulo tempo, anchor=\ancora, inner xsep=2.5pt]
              at (\eixo, -\nivel) {$\nivel$};
          \fi
        }

        \arvoreciclodevida{\instante}
      \end{scope}
    }
  \end{tikzpicture}
  \caption{Ciclo de vida do modelo. A cada passo, os indivíduos da geração
    anterior morrem e seus filhos nascem. Os vértices da geração viva e as
    arestas que chegam até eles aparecem em tom escuro. Os vértices
    transparentes representam a história da evolução da população.}
  \label{fig:ciclo-de-vida}
\end{figure}

\begin{samepage}
Os filhos que um indivíduo deixa formam sua \emph{família}. O tamanho da
família é o número de filhos, que pode ser zero. A árvore genealógica
registra quem é filho de quem, como mostra a \Cref{fig:familia-ramificacao}.
\par
\end{samepage}

\begin{figure}[H]
  \centering
  \begin{tikzpicture}[x=1cm, y=0.72cm]
    \arvoreciclodevida{5}
    \begin{scope}[on background layer]
      \node[fit=(g5b)(g5c), draw=teal, draw opacity=0.65,
            fill=teal, fill opacity=0.10, line width=0.6pt,
            rounded corners=1.5pt, inner sep=2mm] {};
    \end{scope}
  \end{tikzpicture}
  \caption{Uma família no quadro $n=5$ da \Cref{fig:ciclo-de-vida}.
    O retângulo destaca os dois filhos de um mesmo indivíduo.}
  \label{fig:familia-ramificacao}
\end{figure}

Para modelar a reprodução, supomos que:
\begin{enumerate}[label=\Roman*.]
  \item os tamanhos de família são variáveis aleatórias independentes,
        cada uma com valores em $\{0,1,2,\ldots\}$;
  \item os tamanhos de família são variáveis aleatórias identicamente
        distribuídas, com função de probabilidade conhecida, de modo
        que o número $C$ de filhos de um indivíduo típico satisfaz
        $\mathbb{P}(C=k)=p_k$ para $k=0,1,2,\ldots$.
\end{enumerate}
A condição I expressa que cada indivíduo se reproduz sem interferência dos
demais, e a condição II expressa que todos seguem a mesma lei de
probabilidade: o mecanismo de reprodução não muda de um indivíduo para
outro nem de uma geração para outra.

\medskip 

Um processo desse tipo é chamado de \emph{processo de ramificação} e pode
ser usado como modelo simples para fenômenos como o crescimento de
bactérias ou a propagação de um sobrenome de família. Estabelecido o
modelo, o problema básico é descrever como o desenvolvimento do processo
depende da função de probabilidade $p$ dos tamanhos de cada família. 
Para evitar trivialidades, vamos excluir os casos
determinísticos, em que todo indivíduo tem exatamente $k$ filhos com
probabilidade um, e supor ao longo destas notas que
\[
p_k\neq 1
\quad\text{para } k=0,1,2,\ldots.
\]
\begin{figure}[htbp]
  \centering
  \begin{tikzpicture}[x=1cm, y=1.26cm, >=Stealth]
    \draw[eixo tempo, ->] (-0.65, 0.3) -- ([yshift=-2mm]-0.65, -4.25);
    \foreach \geracao/\tamanho in {0/1, 1/3, 2/1, 3/4, 4/7} {
      \draw[guia tempo] (-0.65, -\geracao) -- (10.1, -\geracao);
      \draw[eixo tempo] (-0.65, -\geracao) -- (-0.5, -\geracao);
      \node[rotulo tempo, anchor=east, inner xsep=2.5pt]
        at (-0.65, -\geracao) {$\geracao$};
      \node[font=\small, text=text, anchor=west, inner xsep=2.5pt]
        at (10.1, -\geracao) {$Z_{\geracao}=\tamanho$};
    }

    \foreach \nome/\pai/\geracao/\posicao in {
      g0/raiz/0/7.2,
      g1a/g0/1/4.8, g1b/g0/1/7.2, g1c/g0/1/9.6,
      g2a/g1a/2/4.8,
      g3a/g2a/3/1.2, g3b/g2a/3/3.6, g3c/g2a/3/6, g3d/g2a/3/8.4,
      g4a/g3a/4/0, g4b/g3a/4/1.2, g4c/g3a/4/2.4,
      g4d/g3c/4/5.4, g4e/g3c/4/6.6, g4f/g3d/4/7.8, g4g/g3d/4/9
    } {
      \node[vertice ramificacao] (\nome) at (\posicao, -\geracao) {};
      \ifnum\geracao>0\relax
        \draw[aresta ramificacao] (\pai) -- (\nome);
      \fi
    }
  \end{tikzpicture}
  \caption{Árvore genealógica de um processo de ramificação. Cada ponto
    representa um indivíduo, e cada aresta liga um pai a um de seus filhos.
    As guias tracejadas alinham as gerações ao eixo do tempo, à esquerda,
    e aos respectivos tamanhos $Z_n$, à direita.}
  \label{fig:arvore-ramificacao}
\end{figure}

Chamamos de $n$-ésima geração o conjunto de indivíduos nascidos no tempo
$n$. Denotamos seu tamanho por $Z_n$, com $Z_0=1$. Na
\Cref{fig:arvore-ramificacao}, os valores de $Z_n$ aparecem à direita da árvore.

Os tamanhos das gerações seguintes dependem dos números aleatórios de
filhos. Vamos estudar a sequência $Z_0,Z_1,\ldots$ a partir das seguintes
perguntas:
\begin{enumerate}[label=(\alph*),beginpenalty=10000]
  \item Quais são a média e a variância de $Z_n$?
  \item Qual é a função de probabilidade de $Z_n$?
  \item Qual é a probabilidade de que a população esteja extinta até o
        tempo $n$?
  \item Qual é a probabilidade de que a população se extinga em algum
        momento?
\end{enumerate}

\section{O Método das Funções Geradoras}\label{sec:metodo-geradoras}

O primeiro passo no estudo deste processo de ramificação é explicar como
obter as distribuições das variáveis $Z_n$ em termos da função de
probabilidade $p$ dos tamanhos de cada família. Claramente, $Z_0=1$ e
% Origem: G&W, eq. (9.2).
\begin{equation}\label{eq:zn-um}
\mathbb{P}(Z_1=k)=p_k
\quad\text{para } k=0,1,2,\ldots,
\end{equation}
pois $Z_1$ é o número de filhos da geração inicial.

Para obter a função de probabilidade de $Z_2$, precisamos levar em conta
duas fontes de aleatoriedade. Há $Z_1$ indivíduos na primeira geração, e
cada um deixa um número aleatório de filhos. Escrevendo $C_i$ para o
número de filhos do $i$-ésimo desses indivíduos, temos
\[
Z_2=C_1+C_2+\cdots+C_{Z_1},
\]
uma soma em que tanto o número de parcelas $Z_1$ quanto seus valores $C_i$
são aleatórios, como ilustra a \Cref{fig:soma-aleatoria-z2}.

\begin{figure}[H]
  \centering
  \begin{tikzpicture}[
    x=1cm, y=1.15cm, >=Stealth,
    contorno geracao/.style={caixa familia, draw=red!65, fill=red!9},
  ]
    \draw[eixo tempo, ->] (-2.55, 0.25) -- ([yshift=-2mm]-2.55, -2.25);
    \foreach \geracao in {0,1,2} {
      \draw[eixo tempo] (-2.55, -\geracao) -- (-2.4, -\geracao);
      \node[rotulo tempo, anchor=east, inner xsep=2.5pt]
        at (-2.55, -\geracao) {$n=\geracao$};
    }

    \def\distanciafamilias{1.7}
    \foreach \coluna/\familias/\tamanhos/\parcelas/\soma in {
      0/3/{2,0,3}/{C_1+C_2+C_{Z_1}}/{2+0+3},
      7.6/2/{1,3}/{C_1+C_{Z_1}}/{1+3}
    } {
      \begin{scope}[xshift=\coluna cm, name prefix=quadro\coluna-]
        \pgfmathsetmacro{\meialargura}{0.5*\distanciafamilias*(\familias-1)+0.48}
        \pgfmathsetmacro{\inicioguia}{ifthenelse(\coluna==0,-2.55,-\meialargura-0.37)}
        \foreach \geracao in {0,1,2}
          \draw[guia tempo] (\inicioguia, -\geracao) -- ({\meialargura+0.10}, -\geracao);
        \draw[contorno geracao] (-\meialargura, -1.2) rectangle (\meialargura, -0.8);

        \node[font=\small, text=red!65!black, anchor=west, inner sep=0pt]
          at ({\meialargura+0.23}, -1) {$Z_1=\familias$};
        \node[vertice ramificacao] (raiz) at (0, 0) {};
        \foreach \tamanho [count=\indice] in \tamanhos {
          \pgfmathsetmacro{\centro}{\distanciafamilias*(\indice-(\familias+1)/2)}
          \node[vertice ramificacao] (pai\indice) at (\centro, -1) {};
          \draw[aresta ramificacao] (raiz) -- (pai\indice);
          \node[rotulo tempo, anchor=west, inner xsep=2pt]
            at (pai\indice.east) {$\indice$};

          \ifnum\tamanho=0\relax
            \draw[caixa familia, dashed, fill=white]
              ({\centro-0.58}, -1.8) rectangle ({\centro+0.58}, -2.2);
            \node[rotulo tempo] at (\centro, -2) {$\varnothing$};
          \else
            \draw[caixa familia]
              ({\centro-0.58}, -1.8) rectangle ({\centro+0.58}, -2.2);
            \foreach \filho in {1,...,\tamanho} {
              \node[vertice ramificacao] (filho\indice-\filho)
                at ({\centro+0.35*(\filho-(\tamanho+1)/2)}, -2) {};
              \draw[aresta ramificacao] (pai\indice) -- (filho\indice-\filho);
            }
          \fi
          \node[font=\scriptsize, text=teal!80!black]
            at (\centro, -2.43) {$C_{\indice}=\tamanho$};
        }
        \node[font=\small, text=text, anchor=north, inner sep=0pt] at (0, -2.72)
          {$\begin{aligned}
            Z_2 &= \parcelas \\
                &= \soma
          \end{aligned}$};
      \end{scope}
    }
  \end{tikzpicture}
  \caption{Duas realizações da soma que define $Z_2$. O contorno vermelho
    reúne os $Z_1$ indivíduos da primeira geração. Cada retângulo verde
    contém os $C_i$ filhos do indivíduo $i$ da primeira geração. O retângulo tracejado indica uma
    família vazia. Os exemplos supõem
    \mbox{$p_0,p_1,p_2,p_3>0$}.}
  \label{fig:soma-aleatoria-z2}
\end{figure}

Dado um valor realizado $Z_1=m>0$, as $m$ parcelas $C_1,\ldots,C_m$ são
independentes e têm a mesma função de probabilidade $p$.
Se $Z_1=0$, a soma é vazia e $Z_2=0$.

\bigskip

À primeira vista, parece que podemos descrever o processo
$(Z_n)_{n\geqslant 0}$ considerando apenas uma sequência $C_1,C_2,\ldots$
de variáveis aleatórias independentes, com a mesma função de probabilidade
$p$, para o tamanho de cada família, e escrever, para cada $n$,
\[
Z_n \overset{d}{=} C_1+C_2+\cdots+C_{Z_{n-1}}.
\]
Na igualdade acima, os $C_i$ são cópias genéricas, externas ao processo, e
não os da primeira geração; o símbolo $\overset{d}{=}$ lembra que a
igualdade vale em distribuição, e não entre as variáveis do processo. Como
as funções geradoras de probabilidade dependem apenas das leis, essa versão
forneceria resultados corretos para os cálculos desta seção. Entretanto,
essa tentativa de definir o processo não leva em conta o acoplamento entre
as gerações. No estudo da extinção usaremos que $Z_n=0$ força $Z_{n+1}=0$,
uma propriedade dos eventos definidos pelo processo, e não apenas das leis
de suas variáveis.

Preferimos, por isso, registrar os tamanhos de família que o modelo atribui
a cada indivíduo, com dois índices. Para $n=1,2,\ldots$, temos
\begin{equation}\label{eq:recursao-geracao}
Z_n=C_{n,1}+C_{n,2}+\cdots+C_{n,Z_{n-1}},
\end{equation}
em que $C_{n,k}$ é o número de filhos do $k$-ésimo indivíduo da geração
$n-1$. A \Cref{fig:recursao-geracoes} mostra dois passos da recorrência
na mesma realização do processo.

\begin{figure}[H]
  \centering
  \begin{tikzpicture}[x=1cm, y=3.65cm, >=Stealth]
    \draw[eixo tempo, ->] (-4.7, 3mm) -- ([yshift=-4.5mm]-4.7, -2);
    \foreach \nivel/\tempo in {0/{n-1}, 1/n, 2/{n+1}} {
      \draw[guia tempo] (-4.7, -\nivel) -- (5, -\nivel);
      \draw[eixo tempo] (-4.7, -\nivel) -- (-4.55, -\nivel);
      \node[rotulo tempo, anchor=east, inner xsep=2.5pt]
        at (-4.7, -\nivel) {$\tempo$};
    }

    % As caixas delimitam famílias na geração dos filhos.
    \foreach \nivel/\tempo/\indice/\tamanho/\esquerda/\direita/\rotulo in {
      1/n/1/2/-4.06/-1.65/-3,
      1/n/2/0/-0.54/0.54/0,
      1/n/3/3/0.79/5.03/1.81,
      2/{n+1}/1/1/-4.39/-3.31/-3.85,
      2/{n+1}/2/0/-2.69/-1.61/-2.15,
      2/{n+1}/3/0/0.46/1.54/1,
      2/{n+1}/4/2/2.16/3.36/2.76,
      2/{n+1}/5/0/3.98/5.06/4.52
    } {
      \ifnum\tamanho=0\relax
        \draw[caixa familia, dashed, fill=white]
          ([yshift=-2mm]\esquerda, -\nivel) rectangle
          ([yshift=2mm]\direita, -\nivel);
        \node[rotulo tempo] at (\rotulo, -\nivel) {$\varnothing$};
      \else
        \draw[caixa familia]
          ([yshift=-2mm]\esquerda, -\nivel) rectangle
          ([yshift=2mm]\direita, -\nivel);
      \fi
      \node[font=\scriptsize, text=teal!80!black, yshift=-4.3mm]
        at (\rotulo, -\nivel) {$C_{\tempo,\indice}=\tamanho$};
    }

    \foreach \nome/\posicao/\nivel/\indice in {
      a1/-3/0/1, a2/0/0/2, a3/2.76/0/3,
      b1/-3.85/1/1, b2/-2.15/1/2, b3/1/1/3, b4/2.76/1/4, b5/4.52/1/5,
      c1/-3.85/2/0, c2/2.41/2/0, c3/3.11/2/0
    } {
      \node[vertice ramificacao] (\nome) at (\posicao, -\nivel) {};
      \ifnum\nivel<2\relax
        \node[rotulo tempo, anchor=west, inner xsep=2pt]
          at (\nome.east) {$\indice$};
      \fi
    }
    \foreach \pai/\filho in {a1/b1, a1/b2, a3/b3, a3/b4, a3/b5, b1/c1, b4/c2, b4/c3}
      \draw[aresta ramificacao] (\pai) -- (\filho);

    \node[font=\small, text=red!65!black, anchor=west, inner sep=0pt]
      at (6.28, 0) {$Z_{n-1}=3$};
    \node[font=\small, text=text, anchor=base west, inner sep=0pt]
      at (6.28, -1) {$\begin{aligned}[t]
        Z_n &= \sum\nolimits_{k=1}^{Z_{n-1}} C_{n,k} \\
            &= 2+0+3=5
      \end{aligned}$};
    \node[font=\small, text=text, anchor=base west, inner sep=0pt]
      at (6.28, -2) {$\begin{aligned}[t]
        Z_{n+1} &= \sum\nolimits_{k=1}^{Z_n} C_{n+1,k} \\
                &= 1+0+0+2+0
      \end{aligned}$};
  \end{tikzpicture}
  \caption{Dois passos de uma mesma realização. Em $C_{r,k}$, o índice $r$
    indica a geração em que esta família aparece e $k$ identifica 
    o indivíduo que é o pai desta família.}
  \label{fig:recursao-geracoes}
\end{figure}

Na \Cref{fig:recursao-geracoes}, os três indivíduos da geração $n-1$
deixam $2$, $0$ e $3$ filhos. A soma $2+0+3=5$ conta os cinco indivíduos
que aparecem na linha $n$ da mesma árvore. A soma seguinte tem uma
parcela $C_{n+1,k}$ para cada um desses cinco indivíduos. Nos dois passos,
os tamanhos de cada família seguem a mesma lei $p$.

Na \Cref{fig:soma-aleatoria-z2}, usamos $C_k$ para o número de filhos do
$k$-ésimo indivíduo da primeira geração. Como esses filhos nascem na
geração $2$, escrevemos $C_{2,k}$ para essa mesma variável na notação com
dois índices.


\bigskip 

Somas de um número aleatório de variáveis aleatórias são mais bem
tratadas por funções geradoras de probabilidade do que por funções de
probabilidade. Escrevemos
\begin{equation}\label{eq:def-fgp-geracao}
G_n(s)=\mathbb{E}[s^{Z_n}]
=\sum_{k=0}^{\infty}s^k\,\mathbb{P}(Z_n=k),
\end{equation}
para a função geradora de probabilidade de $Z_n$, e
\begin{equation}\label{eq:def-fgp-familia}
G(s)=\sum_{k=0}^{\infty}s^k\,p_k,
\end{equation}
para a função geradora de probabilidade do tamanho de cada família.
Como $Z_n$ e os tamanhos de família são variáveis aleatórias inteiras e
não negativas, ambas as séries convergem absolutamente para
$|s|\leqslant 1$, pelo argumento apresentado na \Cref{sec:preliminares-fgp}.
Nosso objetivo é expressar $G_n$ em termos de $G$, e fazemos isso no
teorema a seguir.

% Origem: G&W, teorema 9.4.
\begin{theorem}[Iteradas da função geradora]\label{thm:fgp-iterada}
Para $|s|\leqslant 1$, as funções geradoras $G,G_0,G_1,\ldots$ satisfazem
% Origem: G&W, eq. (9.5).
\begin{equation}\label{eq:recursao-fgp}
G_0(s)=s,
\qquad
G_n(s)=G_{n-1}(G(s)),
\quad\text{para } n=1,2,\ldots,
\end{equation}
e, portanto, $G_n$ é a $n$-ésima iterada de $G$:
% Origem: G&W, eq. (9.6).
\begin{equation}\label{eq:fgp-iterada}
G_n(s)=G(G(\cdots G(s)\cdots)),
\quad\text{para } n=0,1,2,\ldots.
\end{equation}
\end{theorem}

\begin{proof}
Como $Z_0=1$ em toda realização, a definição \eqref{eq:def-fgp-geracao}
fornece imediatamente a condição inicial em \eqref{eq:recursao-fgp}:
\begin{equation}\label{eq:fgp-geracao-zero}
G_0(s)=\mathbb{E}[s^{Z_0}]=s.
\end{equation}

Para obter a recorrência em \eqref{eq:recursao-fgp} e, em seguida, a
fórmula iterada \eqref{eq:fgp-iterada}, dividimos a prova para $n\geqslant 1$
em três etapas: primeiro organizamos os sorteios do modelo de modo a
verificar as hipóteses da fórmula da soma aleatória; depois aplicamos o
\Cref{thm:soma-aleatoria} para relacionar $G_n$ a $G_{n-1}$; por fim,
iteramos essa relação por indução.

\smallskip
\noindent\textbf{Etapa 1: Independência entre $Z_{n-1}$ e os tamanhos de
família da geração seguinte.}
Fixe $n\geqslant 1$. Queremos aplicar o \Cref{thm:soma-aleatoria} à
recorrência \eqref{eq:recursao-geracao},
$Z_n=C_{n,1}+C_{n,2}+\cdots+C_{n,Z_{n-1}}$ (com soma vazia igual a $0$
quando $Z_{n-1}=0$), tomando $N=Z_{n-1}$ e $X_k=C_{n,k}$. Para isso, o
teorema exige uma sequência \emph{infinita} $X_1,X_2,\ldots$ de
variáveis independentes e identicamente distribuídas que seja
independente de $N$. Se olharmos apenas para os indivíduos que de fato
nascem na geração $n-1$, o próprio número de variáveis $C_{n,k}$
definidas varia com $Z_{n-1}$, o que tornaria pouco claro até mesmo o
significado de $C_{n,4}$ quando $Z_{n-1}=3$.

Para superar essa dificuldade, representamos todos os sorteios do
modelo de antemão por uma tabela infinita $(C_{r,k})_{r,k\geqslant 1}$
de variáveis aleatórias independentes, todas com função de probabilidade
$p$ e, portanto, com função geradora $G$. A linha $r$ reúne os tamanhos
de família reservados aos eventuais indivíduos da geração $r-1$: sempre
que $Z_{r-1}=m$, a recorrência \eqref{eq:recursao-geracao} consome as $m$
primeiras entradas $C_{r,1},\ldots,C_{r,m}$ da linha $r$ (nenhuma se
$m=0$, caso em que a soma vazia vale $0$), deixando
$C_{r,m+1},C_{r,m+2},\ldots$ sem uso nessa realização. Como as entradas
da tabela são independentes e têm lei $p$, essa construção realiza
exatamente as hipóteses~I e~II do modelo.

Essa representação torna transparente de que dados depende $Z_{n-1}$.
Quando $n=1$, temos $Z_0=1$ constante, logo independente da linha~$1$.
Quando $n\geqslant 2$, podemos acompanhar a evolução linha a linha:
$Z_1=C_{1,1}$ depende apenas da linha~$1$; conhecido $Z_1$, a soma
$Z_2=\sum_{k=1}^{Z_1}C_{2,k}$ envolve apenas as linhas~$1$ e~$2$; de modo
geral, se $Z_{r-1}$ já está determinado pelas linhas $1,\ldots,r-1$,
acrescentar a linha $r$ determina $Z_r$ por \eqref{eq:recursao-geracao}.
Por indução, $Z_{n-1}$ é uma função exclusivamente das linhas
$1,\ldots,n-1$, isto é, de $(C_{r,k})_{1\leqslant r<n,\ k\geqslant 1}$.
Pela independência entre todas as entradas da tabela, a linha $n$
inteira, $(C_{n,k})_{k\geqslant 1}$, é independente da coleção das linhas
anteriores e, consequentemente, de qualquer função delas. Obtemos assim a
independência entre $Z_{n-1}$ e a sequência infinita
$(C_{n,k})_{k\geqslant 1}$, como ilustra a
\Cref{fig:independencia-geracoes}.

\begin{figure}[H]
  \centering
  \begin{tikzpicture}[x=1cm, y=1.4cm, >=Stealth]
    \draw[eixo tempo, ->] (-2.56, 0.2) -- (-2.56, -2.52);
    \foreach \nivel/\tempo in {0/{n-2}, 1/{n-1}, 2/n} {
      \draw[guia tempo] (-2.56, -\nivel) -- (2.48, -\nivel);
      \draw[eixo tempo] (-2.56, -\nivel) -- (-2.41, -\nivel);
      \node[rotulo tempo, anchor=east, inner xsep=2.5pt]
        at (-2.56, -\nivel) {$\tempo$};
    }

    \draw[caixa familia, draw=red!65, fill=red!9]
      (-1.98, -1.14) rectangle (2.13, -0.86);
    \draw[caixa familia] (-2.15, -2.14) rectangle (-1.15, -1.86);
    \draw[caixa familia, dashed, fill=white]
      (-0.5, -2.14) rectangle (0.5, -1.86);
    \node[rotulo tempo] at (0, -2) {$\varnothing$};
    \draw[caixa familia] (1.02, -2.14) rectangle (2.28, -1.86);

    \foreach \nome/\posicao/\nivel in {
      a/0/0, b1/-1.65/1, b2/0/1, b3/1.65/1,
      c1/-1.85/2, c2/-1.45/2, c3/1.25/2, c4/1.65/2, c5/2.05/2
    }
      \node[vertice ramificacao] (\nome) at (\posicao, -\nivel) {};
    \foreach \indice in {1,2,3}
      \node[rotulo tempo, anchor=west, inner xsep=2pt]
        at (b\indice.east) {$\indice$};
    \foreach \pai/\filho in {a/b1, a/b2, a/b3, b1/c1, b1/c2, b3/c3, b3/c4, b3/c5}
      \draw[aresta ramificacao] (\pai) -- (\filho);

    \foreach \centro/\indice/\tamanho in {-1.65/1/2, 0/2/0, 1.65/3/3}
      \node[font=\scriptsize, text=teal!80!black]
        at (\centro, -2.32) {$C_{n,\indice}=\tamanho$};

    \node[font=\small, text=text, anchor=west, inner sep=0pt, align=left]
      at (3, 0) {Reprodução das gerações $0,\ldots,n-2$\\[2pt]
                 $\bigl(C_{r,k}\bigr)_{1\leqslant r<n,\ k\geqslant1}$};
    \draw[eixo tempo, ->] (3.22, -0.34) -- (3.22, -0.79)
      node[midway, right=1mm, font=\scriptsize, text=muted] {determina};
    \node[font=\small, text=red!65!black, anchor=west, inner sep=0pt]
      at (3, -1) {$Z_{n-1}=3$};

    \draw[draw=muted!60, densely dashed, line width=0.45pt]
      (2.92, -1.33) -- (8.45, -1.33);

    \node[font=\small, text=text, anchor=west, inner sep=0pt]
      at (3, -1.65) {Linha $n$, independente das anteriores:};
    \node[font=\small, text=text, anchor=base west, inner sep=0pt]
      at (3, -2)
      {$\underbrace{C_{n,1},\,C_{n,2},\,C_{n,3}}_{\text{parcelas usadas}},\,
        C_{n,4},\ldots$};
  \end{tikzpicture}
  \caption{Separação entre as linhas anteriores e a linha $n$ da tabela.
    As entradas com primeiro índice $r<n$ determinam $Z_{n-1}$ (contorno
    vermelho; aqui $Z_{n-1}=3$). Pela independência das entradas da
    tabela, a linha $n$ inteira, $(C_{n,k})_{k\geqslant 1}$, é independente
    desses dados; nesta realização, entram na soma apenas os três primeiros
    termos ($C_{n,1}=2$, $C_{n,2}=0$, $C_{n,3}=3$), enquanto
    $C_{n,4},\ldots$ ficam sem uso.}
  \label{fig:independencia-geracoes}
\end{figure}

\smallskip
\noindent\textbf{Etapa 2: A recorrência $G_n(s)=G_{n-1}(G(s))$.}
Com as hipóteses verificadas na Etapa~1, aplicamos diretamente o
\Cref{thm:soma-aleatoria} à recorrência \eqref{eq:recursao-geracao},
tomando $N=Z_{n-1}$, $X_k=C_{n,k}$ e $S_N=Z_n$. Como cada parcela
$C_{n,k}$ tem função geradora $G$, a fórmula da soma aleatória
\eqref{eq:fgp-soma-aleatoria} fornece, para todo $|s|\leqslant 1$,
% Origem: G&W, eq. (9.7).
\begin{equation}\label{eq:passo-recursao}
G_n(s)=G_{n-1}\bigl(G(s)\bigr),
\end{equation}
estabelecendo a segunda igualdade de \eqref{eq:recursao-fgp}.

\medskip
\noindent\textbf{Etapa 3: Iteração de $G$.}
Partindo de \eqref{eq:fgp-geracao-zero}, as primeiras aplicações da
recorrência \eqref{eq:passo-recursao} revelam o padrão de composição:
\begin{equation}\label{eq:fgp-primeiras-iteradas}
\begin{aligned}
G_1(s)&=G_0\bigl(G(s)\bigr)=G(s),\\
G_2(s)&=G_1\bigl(G(s)\bigr)=G\bigl(G(s)\bigr),\\
G_3(s)&=G_2\bigl(G(s)\bigr)=G\bigl(G(G(s))\bigr).
\end{aligned}
\end{equation}
Para provar \eqref{eq:fgp-iterada} por indução em $n\geqslant 0$,
notamos primeiro por que podemos seguir aplicando a hipótese de indução
ao ponto $G(s)$: pela estimativa \eqref{eq:fgp-convergencia}, todo $s$
com $|s|\leqslant 1$ satisfaz
$|G(s)|\leqslant\sum_{k=0}^{\infty}p_k|s|^k\leqslant 1$, de modo que o
argumento $u=G(s)$ permanece sempre em $[-1,1]$.

Os casos $n=0$ (identidade em \eqref{eq:fgp-geracao-zero}, zero
aplicações de $G$) e $n=1$ (em \eqref{eq:fgp-primeiras-iteradas}) já
estão verificados. Suponha agora que, para algum $n\geqslant 2$ e para
todo $|u|\leqslant 1$, a função $G_{n-1}(u)$ seja obtida aplicando $G$
sucessivamente $n-1$ vezes a $u$. Fixado $s$ com $|s|\leqslant 1$, o
número $u=G(s)$ pertence a $[-1,1]$; pela recorrência
\eqref{eq:passo-recursao} e pela hipótese de indução aplicada em $u$, o
valor
\[
G_n(s)=G_{n-1}(u)
\]
resulta de $n-1$ aplicações sucessivas de $G$ a $u=G(s)$, isto é, de $n$
aplicações sucessivas de $G$ a $s$:
\[
G_n(s)=G\bigl(G(\cdots G(s)\cdots)\bigr).
\]
Isso conclui a indução e estabelece \eqref{eq:fgp-iterada} para todo
$n\geqslant 0$ e todo $|s|\leqslant 1$.
\end{proof}
% Origem: G&W, seção 9.3, do Teorema 9.8 ao Exercício 9.12.

O \Cref{thm:fgp-iterada} expressa a função geradora de $Z_n$ como
uma iterada de $G$ e contém a informação necessária para estudar
o desenvolvimento do processo. A primeira consequência dessa representação
é a média do tamanho da $n$-ésima geração, que obtemos a seguir.

\begin{theorem}[Média do tamanho da $n$-ésima geração]\label{thm:media-geracao}
Seja $\mu \equiv \sum_{k=0}^{\infty} k p_k$ o valor médio do tamanho de uma
família. Então $\mu\in (0,+\infty]$ e, para cada $n=0,1,2,\ldots$, o
valor médio de $Z_n$ é dado por
\begin{equation}\label{eq:media-mu-n}
\mathbb{E}[Z_n] = \mu^n,
\end{equation}
adotando-se a convenção $\mu^0 \equiv 1$ (inclusive quando $\mu=+\infty$).
\end{theorem}

\begin{proof}
Começamos mostrando que $\mu\in(0,+\infty]$. Pela suposição geral destas
notas de que $p_k\neq 1$ para todo $k$ (em particular, $p_0\neq 1$),
existe algum tamanho de família positivo com probabilidade estritamente
positiva, isto é, $\mathbb{P}(C>0)=\sum_{k=1}^{\infty}p_k=1-p_0>0$. Como
$kp_k\geqslant p_k$ para todo $k\geqslant 1$, segue que
\[
\mu=\sum_{k=1}^{\infty}kp_k\geqslant\sum_{k=1}^{\infty}p_k=1-p_0>0,
\]
logo $\mu\in(0,+\infty]$.

Passamos agora à demonstração de \eqref{eq:media-mu-n} por indução em $n$,
começando pelos casos $n=0$ e $n=1$. Para $n=0$, como $Z_0=1$ e $\mu^0=1$
por convenção, vale imediatamente $\mathbb{E}[Z_0]=1=\mu^0$. Para $n=1$,
a identidade \eqref{eq:zn-um} mostra que $Z_1$ tem a distribuição $p$ do
tamanho de uma família, logo $\mathbb{E}[Z_1]=\mu=\mu^1$ e $G_1=G$; além
disso, o \Cref{thm:fgp-momentos} garante que vale a
seguinte igualdade $G'(1)=\mathbb{E}[Z_1]=\mu$.

Resta o passo indutivo para $n\geqslant 2$. Suponha que
$\mathbb{E}[Z_{n-1}]=\mu^{n-1}$. Pelo \Cref{thm:fgp-momentos}, podemos
obter a média de $Z_n$ como a derivada lateral
$\mathbb{E}[Z_n]=G_n'(1)=\lim_{s\uparrow 1}G_n'(s)$, e o
\Cref{thm:fgp-iterada} fornece a recorrência $G_n(s)=G_{n-1}(G(s))$ em
\eqref{eq:recursao-fgp} para $|s|\leqslant 1$. Não podemos, porém,
aplicar a regra da cadeia diretamente em $s=1$, pois o
\Cref{thm:fgp-momentos} garante a derivabilidade ordinária de uma função
geradora apenas no interior $(-1,1)$ do intervalo de convergência (e a
derivada lateral em $1$ pode inclusive ser $+\infty$). Em vez disso,
calculamos $G_n'(s)$ para $s\in(0,1)$ e depois passamos ao limite quando
$s\uparrow 1$.

Para legitimar tanto a regra da cadeia em $(0,1)$ quanto essa passagem ao
limite, verificamos duas propriedades de $G$:
\begin{enumerate}[label=(\roman*)]
  \item \emph{$G(s)$ permanece em $[0,1)\subset(-1,1)$ para $s\in[0,1)$.}
        Como $p_k\geqslant 0$, $\sum_{k=0}^{\infty}p_k=1$ e existe ao
        menos um $k\geqslant 1$ com $p_k>0$ (pois
        $\sum_{k=1}^{\infty}p_k=1-p_0>0$), o fato de $s^k<1$ para todo
        $k\geqslant 1$ quando $0\leqslant s<1$ produz a desigualdade
        estrita
        \[
        0\leqslant G(s)=p_0+\sum_{k=1}^{\infty}p_k s^k
        <p_0+\sum_{k=1}^{\infty}p_k=1,
        \qquad 0\leqslant s<1.
        \]
        Como $G$ e $G_{n-1}$ são deriváveis em $(-1,1)$ pelo
        \Cref{thm:fgp-momentos} e $G(s)\in[0,1)\subset(-1,1)$, a regra
        da cadeia ordinária aplica-se à composição $G_n=G_{n-1}\circ G$
        em $(0,1)$ e dá
        \[
        G_n'(s)=G_{n-1}'\bigl(G(s)\bigr)\,G'(s),
        \qquad 0<s<1.
        \]
  \item \emph{$G(s)\uparrow 1$ quando $s\uparrow 1$.} Pela continuidade
        de $G$ em $[0,1]$ (estabelecida na \Cref{sec:preliminares-fgp}
        via \Cref{lem:abel}) e pela \Cref{prop:fgp-basica}, temos
        $\lim_{s\uparrow 1}G(s)=G(1)=1$. Como $G$ é não decrescente em
        $[0,1]$ (pois $p_k\geqslant 0$) e satisfaz $G(s)<1$ em $[0,1)$
        pelo item anterior, segue que $G(s)\uparrow 1$ quando
        $s\uparrow 1$.
\end{enumerate}

Passamos agora ao limite $s\uparrow 1$ na expressão obtida para $G_n'(s)$,
analisando cada fator separadamente. Como $G(s)<1$ para $s\in[0,1)$ e
$G(s)\to 1$ quando $s\uparrow 1$, a substituição $u=G(s)$ satisfaz
$u\uparrow 1$; pelo \Cref{thm:fgp-momentos} e pela hipótese de indução,
o primeiro fator converge para
\[
\lim_{s\uparrow 1}G_{n-1}'\bigl(G(s)\bigr)
=\lim_{u\uparrow 1}G_{n-1}'(u)
=G_{n-1}'(1)
=\mathbb{E}[Z_{n-1}]
=\mu^{n-1}.
\]
Para o segundo fator, o \Cref{thm:fgp-momentos} fornece diretamente
$\lim_{s\uparrow 1}G'(s)=G'(1)=\mu$. Como $\mu\in(0,+\infty]$, tanto $\mu^{n-1}$
quanto $\mu$ pertencem a $(0,+\infty]$; assim não ocorre a
indeterminação $0\cdot(+\infty)$, e o limite do produto é o produto dos
limites no sentido estendido (inclusive quando $\mu=+\infty$, em que
$(+\infty)\cdot(+\infty)=+\infty$). Portanto,
\[
\mathbb{E}[Z_n]
=G_n'(1)
=\lim_{s\uparrow 1}G_{n-1}'\bigl(G(s)\bigr)\,G'(s)
=G_{n-1}'(1)\,G'(1)
=\mu^{n-1}\mu
=\mu^n,
\]
o que conclui a indução.
\end{proof}

A variância de $Z_n$ pode ser derivada de modo semelhante, em termos da
média $\mu$ e da variância $\sigma^2$ da distribuição do tamanho de uma família;
o cálculo usa a segunda derivada de $G_n$ em $1$ e está proposto
no \Cref{exercise:variancia-geracao}.


\bigskip

Além de fornecer o valor médio em cada geração, o \Cref{thm:media-geracao}
dá a primeira indicação sobre o destino da população no longo prazo.
Fazendo $n\to\infty$ em \eqref{eq:media-mu-n}, o comportamento assintótico
da média exibe uma transição brusca em $\mu=1$:
\[
\mathbb{E}[Z_n] \to
\begin{cases}
0 & \text{se } \mu<1,\\
1 & \text{se } \mu=1,\\
\infty & \text{se } \mu>1,
\end{cases}
\]
o que indica que o comportamento do processo depende substancialmente
de qual dos três casos $\mu<1$, $\mu=1$ ou $\mu>1$ ocorre.
Veremos esse contraste com mais detalhe na \Cref{sec:extincao}:
se $\mu\leqslant 1$, a extinção é certa, enquanto se $\mu>1$ há
probabilidade estritamente positiva de que a linhagem continue
indefinidamente.

Essa dependência da média $\mu$ do tamanho de família é natural.
A condição $\mu<1$ significa que cada morte é reposta, em média,
por menos de um nascimento: cada indivíduo gera, em média, estritamente
menos filhos do que o necessário para preencher a lacuna deixada pela sua
morte.
Já $\mu>1$ significa que cada morte é reposta, em média, por mais
de um nascimento, de modo que a população aumenta em média.

O caso $\mu=1$ é chamado \emph{crítico}: a média do tamanho da população
é igual a $1$ em todos os tempos. Isso não significa que a população
permaneça parada em $1$; ela oscila aleatoriamente, e essas flutuações
fazem com que o tamanho da população assuma o valor $0$ em algum momento.
A partir daí, a população está extinta.

% Origem: G&W, exercício 9.10.
\begin{exercise}\label{exercise:composicao-geracoes}
Mostre que, no processo de ramificação acima,
\[
G_n(s) = G_r\bigl(G_{n-r}(s)\bigr),
\]
para todo $r = 0, 1, 2, \ldots, n$. Isso pode ser provado diretamente
a partir da conclusão do \Cref{thm:fgp-iterada} ou adaptando o método
de prova de \eqref{eq:passo-recursao}.
\end{exercise}

% Origem: G&W, exercício 9.11.
\begin{exercise}\label{exercise:extincao-bernoulli}
Suponha que cada família de um processo de ramificação
contenha apenas um membro (com probabilidade $p$) ou seja vazia
(com probabilidade $1-p$). Encontre a probabilidade de que o processo
se extinga até a $n$-ésima geração.
\end{exercise}

% Origem: G&W, exercício 9.12.
\begin{exercise}\label{exercise:variancia-geracao}
Sejam $\mu$ e $\sigma^2$ a média e a variância da distribuição do
tamanho de família. Adapte a prova do \Cref{thm:media-geracao} para
mostrar que a variância de $Z_n$, o tamanho da $n$-ésima geração do
processo de ramificação, é dada por
\[
\operatorname{Var}(Z_n) =
\begin{cases}
n\sigma^2 & \text{se } \mu=1,\\[0.4cm]
\sigma^2\mu^{n-1}\, \displaystyle\frac{\mu^n-1}{\mu-1} & \text{se } \mu\neq 1.
\end{cases}
\]
\end{exercise}
% Origem: G&W, seção 9.4 (primeiro caso).
\section{Extinção no Caso Geométrico}\label{sec:exemplo}

A equação funcional \eqref{eq:recursao-fgp} do \Cref{thm:fgp-iterada},
$G_n=G_{n-1}\circ G$, relaciona as funções geradoras de duas gerações
consecutivas. Em geral, não é possível obter uma expressão explícita
para $G_n$, o que só acontece em poucos casos particulares. Nesta seção,
estudamos um desses casos: embora a distribuição do tamanho de família
seja bastante simples, o modelo correspondente já exibe o fenômeno da
extinção.

Suponhamos que a função de probabilidade do tamanho de uma família seja
\[
p_k=\left(\frac{1}{2}\right)^{k+1}
\quad\text{para } k=0,1,2,\ldots.
\]
A normalização segue da soma da série geométrica,
$\sum_{k\geqslant 0}p_k=\tfrac{1}{2}\sum_{k\geqslant 0}(1/2)^k=1$, de modo
que a família vazia ocorre com probabilidade $1/2$ e os tamanhos pequenos
são os mais prováveis.

Essa lei é, a menos de uma unidade, a distribuição geométrica de parâmetro
$\tfrac{1}{2}$. De fato, se $Y\sim\mathrm{Geom}\bigl(\tfrac{1}{2}\bigr)$, de
modo que $\mathbb{P}(Y=j)=(1/2)^{j}$ para $j\geqslant 1$, então $C=Y-1$
assume os valores $0,1,2,\ldots$ e
\[
\mathbb{P}(C=k)=\mathbb{P}(Y=k+1)=\left(\frac{1}{2}\right)^{k+1}=p_k.
\]
Da igualdade acima segue que o tamanho de uma família pode ser interpretado
como o número de tentativas até o primeiro
sucesso, contando a tentativa vencedora, menos um. Nas notas de aula sobre
funções geradoras de probabilidade calculamos a função geradora da distribuição
geométrica e discutimos esse deslocamento do suporte; o cálculo a seguir repete
aquele, com um fator $s$ a menos, e por isso o apresentamos em detalhe.

Como $p$ decai geometricamente, a função geradora de probabilidade do
tamanho de uma família é um múltiplo de uma série geométrica de razão $s/2$:
\begin{equation}\label{eq:fgp-geometrica-meio}
G(s)=\sum_{k=0}^{\infty}s^k\left(\frac{1}{2}\right)^{k+1}
=\frac{1}{2}\sum_{k=0}^{\infty}\left(\frac{s}{2}\right)^{k}
=\frac{1}{2-s},
\qquad |s|\leqslant 1.
\end{equation}
A série do meio converge para $|s/2|<1$, isto é, para $|s|<2$, intervalo em que
vale a igualdade final; como todas as funções geradoras de probabilidade do
processo são usadas em $|s|\leqslant 1$, trabalhamos nesse domínio, no qual o
denominador $2-s$ não se anula, pois $2-\operatorname{Re}(s)\geqslant 1$. Note ainda que
$G_Y(s)=s/(2-s)=s\,G(s)$: a multiplicação por $s$ é o efeito de acrescentar
uma unidade a $C$, isto é, de passar de $C$ para $Y=C+1$.

Com $G$ em mãos, aplicamos a equação funcional sucessivamente. Para $n=1$ não
há composição a fazer, pois $Z_1$ tem a distribuição do tamanho de uma família:
\[
G_1(s)=G(s)=\frac{1}{2-s}.
\]
Para $n=2$, a composição de $G$ com ela mesma produz uma fração de frações, que
simplificamos em seguida:
\begin{align*}
G_2(s)
&=G\bigl(G_1(s)\bigr)
=\frac{1}{2-G_1(s)}
=\frac{1}{2-\dfrac{1}{2-s}}
\\[0.3cm]
&=\frac{2-s}{2(2-s)-1}
=\frac{2-s}{3-2s}.
\end{align*}
Para $n=3$, o mesmo procedimento, agora aplicado a $G_2$,
dá
\[
G_3(s)
=G_2\bigl(G(s)\bigr)
=\frac{2-G(s)}{3-2G(s)}
=\frac{2-\dfrac{1}{2-s}}{3-\dfrac{2}{2-s}}
=\frac{3-2s}{4-3s}.
\]

Os três resultados seguem o mesmo padrão. Em $G_n$, o numerador é
$n-(n-1)s$ e o denominador é $n+1-ns$: para $n=1$ isso reproduz $1/(2-s)$,
para $n=2$ temos $(2-s)/(3-2s)$ e para $n=3$, $(3-2s)/(4-3s)$. Isso sugere a
seguinte fórmula, válida para todo $n\geqslant 1$ e todo $|s|\leqslant 1$:
\begin{equation}\label{eq:gn-meio}
G_n(s)=\frac{n-(n-1)s}{n+1-ns}.
\end{equation}
O denominador em \eqref{eq:gn-meio} não se anula no domínio considerado, pois
sua parte real é $n+1-n\operatorname{Re}(s)\geqslant 1$, de modo que a
expressão está bem definida em todo $|s|\leqslant 1$.

Vamos demonstrar a fórmula \eqref{eq:gn-meio} por indução em $n$.

\begin{proof}
Para $n=1$, a substituição direta em \eqref{eq:gn-meio} dá $1/(2-s)$, que é o
valor já obtido para $G_1$. Suponhamos agora que a fórmula vale para $n-1$,
com $n\geqslant 2$, e deduzamos o caso $n$. Pela equação funcional
\eqref{eq:recursao-fgp}, escrita como $G_n=G\circ G_{n-1}$,
\[
G_n(s)=G\bigl(G_{n-1}(s)\bigr)=\frac{1}{2-G_{n-1}(s)}.
\]
A composição é legítima porque $|G_{n-1}(s)|\leqslant 1$ para $|s|\leqslant 1$,
conforme observamos na \Cref{sec:preliminares-fgp}, e porque $2-G_{n-1}(s)$ não
se anula. Substituindo a hipótese de indução e reunindo os termos sobre um único
denominador,
\begin{align*}
G_n(s)
&=\frac{1}{2-\dfrac{(n-1)-(n-2)s}{n-(n-1)s}}
\\[0.3cm]
&=\frac{n-(n-1)s}{2\bigl(n-(n-1)s\bigr)-\bigl((n-1)-(n-2)s\bigr)}
\\[0.3cm]
&=\frac{n-(n-1)s}{2n-2(n-1)s-n+1+(n-2)s}
\\[0.3cm]
&=\frac{n-(n-1)s}{n+1-ns}.
\end{align*}
Na última passagem, os termos constantes do denominador produzem
$2n-(n-1)=n+1$, e os coeficientes de $s$ produzem
$-2(n-1)+(n-2)=-2n+2+n-2=-n$, o que é exatamente a fórmula
\eqref{eq:gn-meio} para o índice $n$. Isso encerra a indução.
\end{proof}

A fórmula \eqref{eq:gn-meio} dá acesso à função de probabilidade de $Z_n$,
pois os coeficientes da série de potências de uma função geradora são essas
probabilidades. Para extraí-los, convém decompor o quociente em uma constante
mais uma série geométrica. Escrevemos o numerador como múltiplo do denominador
mais uma constante:
\[
n-(n-1)s=\frac{n-1}{n}\bigl(n+1-ns\bigr)+\frac{1}{n}.
\]
Dividindo por $n+1-ns$ e, em seguida, normalizando o denominador da fração
restante para que o termo constante seja $1$, obtemos
\begin{align*}
G_n(s)
&=\frac{n-1}{n}+\frac{1}{n}\cdot\frac{1}{n+1-ns}
\\[0.3cm]
&=\frac{n-1}{n}+\frac{1}{n(n+1)}\cdot\frac{1}{1-\frac{n}{n+1}\,s}.
\end{align*}
O último fator é a soma da série geométrica de razão
$\tfrac{n}{n+1}s$, cujo módulo é menor do que $1$ quando $|s|\leqslant 1$. O
coeficiente de $s^0$ reúne o termo constante com o primeiro termo da série,
\[
\frac{n-1}{n}+\frac{1}{n(n+1)}
=\frac{(n-1)(n+1)+1}{n(n+1)}
=\frac{n}{n+1},
\]
enquanto, para $k\geqslant 1$, o coeficiente vem apenas da série:
\[
\frac{1}{n(n+1)}\left(\frac{n}{n+1}\right)^{k}
=\frac{n^{k-1}}{(n+1)^{k+1}}.
\]
Reunindo os dois casos,
\begin{equation}\label{eq:distribuicao-zn-meio}
\mathbb{P}(Z_n=k)=
\begin{cases}
\dfrac{n}{n+1}, & k=0,\\[0.4cm]
\dfrac{n^{k-1}}{(n+1)^{k+1}}, & k=1,2,\ldots.
\end{cases}
\end{equation}
Uma verificação de consistência é imediata: os termos com $k\geqslant 1$ somam
$1/(n+1)$, pois formam uma série geométrica de razão $n/(n+1)$ multiplicada por
$1/(n+1)^2$, e a soma com $n/(n+1)$ fornece $1$, como a normalização exige.

O caso $k=0$ em \eqref{eq:distribuicao-zn-meio} é o que descreve a extinção até
a $n$-ésima geração, e sua probabilidade vale
\[
\mathbb{P}(Z_n=0)=\frac{n}{n+1}=1-\frac{1}{n+1}\longrightarrow 1
\quad\text{quando } n\to\infty.
\]

Já que 
\begin{align*}
\{Z_{n}=0\}  \subseteq \{Z_{n+1}=0\},\qquad \forall n\in\mathbb{N}\cup\{0\}
\end{align*}
temos da propriedade de continuidade da medida que 
\begin{align*}
\mathbb{P}\left( \bigcup_{n=0}^{\infty} \{Z_n=0\} \right) 
=
\lim_{n\to\infty} \mathbb{P}(Z_n=0) 
=
\lim_{n\to\infty} 1-\frac{1}{n+1} = 1. 
\end{align*} 
Neste exemplo, portanto, a
extinção ocorre com probabilidade $1$.


\medskip 
É importante observar como os cálculos feito até aqui fornecem 
informações complementares sobre o modelo. Por exemplo, já que 
$G'(s)=1/(2-s)^2$, o valor médio do tamanho de uma
família é $\mu=G'(1)=1$, de modo que o processo está no caso crítico e, pelo
\Cref{thm:media-geracao}, $\mathbb{E}[Z_n]=1$ para todo $n$. A média permanece
constante enquanto $\mathbb{P}(Z_n=0)\to 1$, e isso não é contraditório: no
evento, de probabilidade cada vez menor, em que a população não se extingue,
ela assume valores cada vez maiores, e é essa massa que sustenta a média. Na
\Cref{sec:extincao} confirmamos, para um processo geral, que $\mu\leqslant 1$
implica extinção quase certa.


\bigskip 

% Origem: G&W, seção 9.4 (Um exemplo), caso geral e exercício 9.17.
O exemplo analisado acima é o caso particular em que cada tamanho de uma família tem
distribuição geométrica deslocada, de parâmetro $1/2$. Vale a pena
considerar a família geométrica completa, porque ela permite acompanhar o efeito
da média do tamanho de família sobre a extinção sem abandonar o cálculo
explícito. Suponha então que a função de probabilidade de cada tamanho de
família seja
\[
p_k=pq^k,
\qquad k=0,1,2,\ldots,
\]
onde $0<p=1-q<1$. Neste caso, $p_0=p>0$, de modo que uma família vazia é
possível e a extinção já pode ocorrer na primeira geração. 
Neste caso o valor médio do tamanho de uma família é dado pela seguinte expressão
\[
\mu=\sum_{k=0}^{\infty}k\,pq^k
=pq\sum_{k=1}^{\infty}k\,q^{k-1}
=\frac{pq}{(1-q)^2}
=\frac{q}{p},
\]
onde na segunda igualdade usamos a série geométrica e sua derivada.

A função geradora do tamanho de família é a soma de uma série
geométrica. Mais precisamente, para $|qs|<1$, isto é, para $|s|<q^{-1}$,
\begin{equation}\label{eq:fgp-geometrica-geral}
G(s)=\sum_{k=0}^{\infty}s^k pq^k
=p\sum_{k=0}^{\infty}(qs)^k
=\frac{p}{1-qs}.
\end{equation}
O raio de convergência desta série é $q^{-1}>1$, portanto o domínio de
\eqref{eq:fgp-geometrica-geral} contém o intervalo $|s|\leqslant 1$, que é o
único que usaremos daqui em diante. Quando $p=q=1/2$, a expressão
\eqref{eq:fgp-geometrica-geral} se reduz a $1/(2-s)$, como no exemplo anterior.
Convém supor agora $p\neq q$, ou seja, $\mu\neq 1$, já que o caso $\mu=1$ é exatamente
o exemplo anterior.


\medskip 

Iterando \eqref{eq:fgp-geometrica-geral} obtemos uma fórmula fechada para
$G_n$. Como a composição de funções é associativa e $G_1=G$, a recursão
\eqref{eq:recursao-fgp} pode ser escrita também na forma
$G_n=G\circ G_{n-1}$. Partindo dessa forma, uma indução em $n$ mostra que, para
$n=1,2,\ldots$ e $|s|\leqslant 1$,
\begin{equation}\label{eq:gn-geral}
G_n(s)
=p\,\frac{(q^n-p^n)-qs\,(q^{n-1}-p^{n-1})}
{(q^{n+1}-p^{n+1})-qs\,(q^n-p^n)}.
\end{equation}
No passo indutivo basta substituir $G_{n-1}(s)$ na expressão de $G$ em
\eqref{eq:fgp-geometrica-geral} e simplificar as frações que aparecem; a
álgebra é rotineira e não a reproduzimos aqui. Para conferir o caso base, note
que em $n=1$ o numerador de \eqref{eq:gn-geral} vale $p(q-p)$ e o denominador
vale $(q-p)(q+p)-qs(q-p)=(q-p)(1-qs)$, e a razão devolve
$p/(1-qs)=G(s)$.



\medskip 
A probabilidade de extinção até a $n$-ésima geração é o valor de $G_n$ em $0$,
que corresponde a manter apenas o termo constante da iterada. Tomando $s=0$ em
\eqref{eq:gn-geral} e usando $\mu=q/p$, obtemos
\begin{equation}\label{eq:extincao-geometrica}
\mathbb{P}(Z_n=0)
=G_n(0)
=p\,\frac{q^n-p^n}{q^{n+1}-p^{n+1}}
=\frac{\mu^n-1}{\mu^{n+1}-1}.
\end{equation}
A última igualdade resulta de colocar $p^{n+1}$ em evidência no numerador e no
denominador: $p(q^n-p^n)=p^{n+1}(\mu^n-1)$ e
$q^{n+1}-p^{n+1}=p^{n+1}(\mu^{n+1}-1)$. Assim, toda a dependência em $n$ fica
expressa por meio de uma única quantidade, a média $\mu$ do tamanho de família.

O comportamento assintótico de \eqref{eq:extincao-geometrica} depende de $\mu$.
Se $\mu<1$, então $\mu^n\to 0$ e tanto o numerador quanto o denominador tendem a
$-1$, de modo que
\[
\mathbb{P}(Z_n=0)\longrightarrow 1.
\]
Se $\mu>1$, dividimos numerador e denominador por $\mu^{n+1}$ e escrevemos
\[
\mathbb{P}(Z_n=0)
=\frac{\mu^n-1}{\mu^{n+1}-1}
=\frac{\mu^{-1}-\mu^{-n-1}}{1-\mu^{-n-1}}
\longrightarrow \frac{1}{\mu}.
\]
Reunindo os dois casos,
\[
\mathbb{P}(Z_n=0)\longrightarrow
\begin{cases}
1 & \text{se } \mu<1,\\[0.2cm]
\mu^{-1} & \text{se } \mu>1.
\end{cases}
\]

No primeiro caso, a probabilidade de que a população já esteja extinta na
$n$-ésima geração tende a $1$ e, uma vez extinta a população, ela permanece
extinta, de modo que a extinção em algum momento tem probabilidade $1$. 
Mais precisamente, como argumentado no exemplo anterior segue 
de $\{Z_n=0\}\subseteq \{Z_{n+1}=0\}$ que se $\mu<1$, então
\begin{align*}
\mathbb{P}\left( \bigcup_{n=0}^{\infty} \{Z_n=0\} \right) 
=
\lim_{n\to\infty} \mathbb{P}(Z_n=0) 
=
1. 
\end{align*} 


No segundo caso o mesmo limite é $1/\mu<1$: para todo $n$ suficientemente grande a
probabilidade de extinção até a $n$-ésima geração fica abaixo de um número
estritamente menor que $1$, o número esperado de filhos por indivíduo supera a
reposição, e a extinção em algum momento pode não ocorrer com probabilidade positiva,
ou seja, há portanto probabilidade positiva de que a linhagem sobreviva em todos os tempos.
Acrescentando o caso $p=q=1/2$, em que $\mu=1$ e a extinção em algum momento
também tem probabilidade $1$, como vimos no exemplo anterior, chegamos à
conclusão de que a extinção é certa se, e somente se, $\mu\leqslant 1$. 
Na próxima seção veremos que esse critério não depende da forma geométrica da função de
probabilidade: ele vale para qualquer processo de ramificação com tamanhos de
família não degenerados.

% Origem: G&W, exercício 9.17.
\begin{exercise}\label{exercise:media-variancia-geometrica}
Encontre a média e a variância de $Z_n$ quando a distribuição do tamanho de
família é dada por $p_k=pq^k$ para $k=0,1,2,\ldots$, com $0<p=1-q<1$. Deduza
que $\operatorname{Var}(Z_n)\to 0$ se e somente se $p>1/2$.
\end{exercise}

% Origem: G&W, seção 9.5 (The probability of extinction).
\section{O Critério Geral de Extinção}\label{sec:extincao}

No caso geométrico da seção anterior, o processo de ramificação se extingue
com probabilidade um se, e somente se, a média $\mu$ do tamanho de família
satisfaz $\mu\leqslant 1$. Vamos mostrar nesta seção um resultado que
reduz o cálculo da probabilidade de extinção, em qualquer processo de
ramificação, ao estudo de uma equação envolvendo a função geradora $G$.
A probabilidade de que a população se extinga em algum momento é dada por
\begin{equation}\label{eq:definicao-extincao}
e\equiv \mathbb{P}\bigl(Z_n=0, \ \text{para algum } n\geqslant 0\bigr).
\end{equation}
O evento cuja probabilidade é $e$ descreve a situação em que a linhagem
se encerra em tempo finito. Como $Z_0=1$, esse evento não
pode ocorrer em $n=0$, mas o índice começa em zero porque a geração inicial
faz parte da descrição do processo.

Para calcular $e$, aproximamos esse evento por eventos em que a extinção
já ocorreu até um tempo prefixado. Escrevemos
\[
E_n=\{Z_n=0\}
\]
para o evento de que a população já esteja extinta na $n$-ésima geração e
denotamos a probabilidade deste evento por
\[ 
e_n\equiv \mathbb{P}(E_n).
\] 

A extinção em algum momento é o evento dado pela união de todos os eventos $E_n$'s, isto é,
\[
\bigl\{Z_n=0, \ \text{para algum } n\geqslant 0\bigr\}
=\bigcup_{n=0}^{\infty}E_n.
\]

Como já mencionado anterioremente a família de eventos $(E_n)_{n\in\mathbb{N}}$ 
é uma família crescente de eventos. 
De fato, se $Z_n=0$, então a $n$-ésima
geração é vazia, nenhum indivíduo existe para se reproduzir e a geração
seguinte também é vazia, isto é, $Z_{n+1}=0$. Logo
$E_n\subseteq E_{n+1}$ e, em particular, $e_n\leqslant e_{n+1}$. Para uma
sequência crescente de eventos, a continuidade da probabilidade, discutida
nas notas de aula sobre probabilidade, identifica a probabilidade da
reunião com o limite das probabilidades. Logo temos
\begin{equation}\label{eq:limite-en}
e=\lim_{n\to\infty}e_n .
\end{equation}
A igualdade \eqref{eq:limite-en} transfere o problema para a sequência
$(e_n)_{n\in\mathbb{N}}$, que é mais tratável porque cada $e_n$ é a probabilidade de um
evento descrito em um único instante. Além disso, como a sequência é
crescente, ela tem limite, e esse limite é o valor procurado mesmo quando
as probabilidades $e_n$ não são conhecidas explicitamente.

Um caso extremo é imediato. Se $p_0=0$, todas as famílias são não vazias,
isto é, todo indivíduo tem ao menos um filho. Nesse caso, por indução em $n$,
temos $Z_n\geqslant 1$ para todo $n$: a geração inicial é não vazia e, se
$Z_n\geqslant 1$, então $Z_{n+1}$ é uma soma de $Z_n$ tamanhos de família,
todos eles pelo menos $1$. Assim
$\mathbb{P}(Z_n=0)=0$ para todo $n$ e, por \eqref{eq:limite-en}, $e=0$.
O interesse está, portanto, no caso em que $p_0>0$, e é a ele que se
dedica o resultado a seguir.

\begin{theorem}[Menor raiz da equação de extinção]\label{thm:extincao-raiz}
A probabilidade $e$ de extinção em algum momento é a menor raiz não
negativa da equação
\begin{equation}\label{eq:equacao-extincao}
x=G(x).
\end{equation}
\end{theorem}

\begin{proof}
Primeiro identificamos $e_n$ com um valor da função geradora de $Z_n$.
Como $e_n=\mathbb{P}(Z_n=0)$ e $G_n$ é a função geradora de $Z_n$, a
\Cref{prop:fgp-basica} aplicada a $Z_n$ fornece
\[
e_n=G_n(0)
\quad\text{para } n=0,1,2,\ldots .
\]
Precisamos agora de uma relação de recorrência entre os $e_n$. Para
$|s|\leqslant 1$, o \Cref{thm:fgp-iterada} garante que $G_n$ é a
$n$-ésima iterada de $G$. Como a composição de funções é associativa,
podemos agrupar essa iterada de duas maneiras, e obtemos
\[
G_n(s)=G\bigl(G_{n-1}(s)\bigr),
\]
que é a identidade \eqref{eq:recursao-fgp} lida na ordem oposta, e também
o caso particular de \eqref{eq:fgp-iterada} em que a primeira das $n$
aplicações de $G$ é destacada. Avaliando em $s=0$, obtemos
\begin{equation}\label{eq:recursao-en}
e_n=G(e_{n-1})
\quad\text{para } n=1,2,\ldots ,
\end{equation}
com a condição inicial $e_0=G_0(0)=0$.

Passamos ao limite em \eqref{eq:recursao-en}. O lado esquerdo tende a $e$
por \eqref{eq:limite-en}. Para o lado direito, observamos que $G$ é uma
série de potências com raio de convergência pelo menos $1$ e coeficientes
não negativos; como vimos na \Cref{sec:preliminares-fgp}, isso faz de $G$
uma função contínua em $[0,1]$. Como todos os $e_n$ pertencem a $[0,1]$,
podemos trocar o limite pela função:
\[
e=\lim_{n\to\infty}e_n
=\lim_{n\to\infty}G(e_{n-1})
=G\Bigl(\lim_{n\to\infty}e_{n-1}\Bigr)
=G(e).
\]
Logo $e$ é uma raiz não negativa de \eqref{eq:equacao-extincao}.

Resta mostrar que $e$ não é maior que qualquer outra raiz não negativa.
Seja então $\eta\geqslant 0$ uma solução qualquer de \eqref{eq:equacao-extincao},
de modo que $G(\eta)=\eta$. Se $\eta>1$, então $e\leqslant 1<\eta$ e não há
nada a provar; suponha portanto $0\leqslant\eta\leqslant 1$. Como os
coeficientes $p_k$ de $G$ são não negativos, $G$ é não decrescente em
$[0,1]$: se $0\leqslant u\leqslant v\leqslant 1$,
então $u^k\leqslant v^k$ para todo $k$ e, multiplicando por $p_k\geqslant 0$
e somando, $G(u)\leqslant G(v)$. A partir da condição inicial, a
recorrência \eqref{eq:recursao-en} produz
\[
e_1=G(e_0)=G(0)\leqslant G(\eta)=\eta .
\]
Repetindo o argumento, e usando agora $e_1\leqslant\eta$,
\[
e_2=G(e_1)\leqslant G(\eta)=\eta .
\]
Por indução em $n$, obtemos $e_n\leqslant\eta$ para todo $n\geqslant 1$.
Fazendo $n\to\infty$ e usando \eqref{eq:limite-en}, concluímos que
$e\leqslant\eta$.

Assim $e$ é uma raiz não negativa de \eqref{eq:equacao-extincao} e é menor
ou igual a todas as demais, o que encerra a demonstração.
\end{proof}

O \Cref{thm:extincao-raiz} troca um problema sobre a sequência de eventos
$E_n$ por um problema sobre os pontos fixos de $G$. A equação
\eqref{eq:equacao-extincao} sempre tem a solução $x=1$, pois $G(1)=1$ pela
\Cref{prop:fgp-basica}, de modo que a probabilidade de extinção é o menor
ponto fixo de $G$ no intervalo $[0,1]$. Em particular, $e=1$ exatamente
quando $1$ é o único ponto fixo de $G$ nesse intervalo.

\begin{exercise}\label{exercise:ponto-fixo-intervalo}
Mostre que, se $f:[0,1]\to [0,1]$ é uma função contínua, então $f$ possui
pelo menos um ponto fixo, isto é, existe algum $x_0\in [0,1]$ tal que
$f(x_0)=x_0$. Para isso, considere a função auxiliar $g(x)=f(x)-x$ e aplique o
Teorema do Valor Intermediário.
\end{exercise}


\bigskip 

A descrição apresentada acima é exata, mas por si só não decide a questão 
que mais interessa na prática, isto é,
se a extinção é certa ou se há probabilidade positiva de que a população
sobreviva indefinidamente. Veremos que a resposta depende
apenas do valor médio $\mu$ do tamanho de uma família e que a fronteira entre os dois
comportamentos ocorre em $\mu=1$.

\begin{theorem}[Critério de extinção]\label{thm:extincao-criterio}
Suponha que $p_1\neq1$. Então
\[
e=1
\quad \Longleftrightarrow \quad
\mu\leqslant1,
\]
onde $\mu=\sum_{k}kp_k$ é o valor médio do tamanho de uma família e o número $e$ é a
probabilidade de extinção definida em \eqref{eq:definicao-extincao}.
\end{theorem}

O caso $p_1=1$ está excluído do enunciado porque nele a conclusão é falsa. Se
$p_1=1$, cada família tem tamanho $1$ com probabilidade um e,
consequentemente, $Z_n=1$ para todo $n$: a população nunca se extingue, de modo
que $e=0$, enquanto $\mu=1$. A hipótese $p_1\neq1$ é mais fraca do que a
suposição geral adotada nestas notas, segundo a qual $p_k\neq1$ para todo $k$,
e o critério vale nesse ambiente mais amplo. Note também que o critério inclui
o caso $p_0=0$, estudado no início da seção, no qual a extinção é impossível e,
como veremos na prova, a média satisfaz $\mu>1$.

Do ponto de vista geométrico, as raízes de $x=G(x)$ são as abscissas dos pontos
em que o gráfico de $G$ encontra a diagonal $y=x$ do quadrado $[0,1]^{2}$. Como
$G$ é não decrescente e convexa, com $G(1)=1$, esse gráfico sempre passa por
$(1,1)$; o que a inclinação de $G$ em $1$ decide é se há ou não um segundo
encontro no interior do quadrado. A \Cref{fig:extincao-raizes} registra os dois
casos. A prova que apresentamos a seguir é analítica e não depende do desenho:
ela compara $G$ com a diagonal por meio da função $h(x)=G(x)-x$ e usa apenas
monotonicidade, convexidade e o teorema do valor intermediário.

\begin{figure}[htbp]
  \centering
  \begin{tikzpicture}[x=5.4cm, y=5.4cm, >=Stealth,
      declare function={Ga(\x)=0.2+0.1*\x+0.7*\x*\x;
                        Gb(\x)=0.2+0.7*\x+0.1*\x*\x;}]
    \tikzset{
      eixo/.style={black!65, line width=0.55pt, ->},
      tick/.style={black!65, line width=0.55pt},
      rotulo/.style={font=\scriptsize, text=muted},
      diagonal/.style={densely dashed, black!45, line width=0.7pt},
      curva/.style={teal, line width=1.35pt},
      raiz/.style={fill=teal!85!black},
      nota/.style={font=\scriptsize, text=teal!85!black},
    }

    % Painel (a): G'(1)=1.5>1. As raízes de G(x)=x em [0,1] são e=2/7 e 1.
    \begin{scope}
      \draw[eixo] (-0.02,0) -- (1.14,0) node[below, rotulo] {$x$};
      \draw[eixo] (0,-0.02) -- (0,1.14) node[left, rotulo] {$y$};
      \foreach \v/\rot in {0.5/{0{,}5},1/1} {
        \draw[tick] (\v,0.013) -- (\v,-0.013) node[below=1pt, rotulo] {$\rot$};
        \draw[tick] (0.013,\v) -- (-0.013,\v) node[left=1pt, rotulo] {$\rot$};
      }
      \node[below left=1pt, rotulo] at (0,0) {$0$};
      \draw[diagonal] (0,0) -- (1.05,1.05);
      \draw[curva] plot[domain=0:1, samples=100, smooth] (\x, {Ga(\x)});
      \fill[raiz] (2/7, 2/7) circle (2pt);
      \fill[raiz] (1, 1) circle (2pt);
      \node[below left=2pt, nota] at (2/7, 2/7) {$e$};
      \node[below right=2pt, nota] at (1, 1) {$1$};
      \node[anchor=south, font=\small, text=text] at (0.57, 1.14) {(a) $G'(1)>1$};
    \end{scope}

    % Painel (b): G'(1)=0.9<=1. A única raiz de G(x)=x em [0,1] é x=1.
    \begin{scope}[xshift=7.6cm]
      \draw[eixo] (-0.02,0) -- (1.14,0) node[below, rotulo] {$x$};
      \draw[eixo] (0,-0.02) -- (0,1.14) node[left, rotulo] {$y$};
      \foreach \v/\rot in {0.5/{0{,}5},1/1} {
        \draw[tick] (\v,0.013) -- (\v,-0.013) node[below=1pt, rotulo] {$\rot$};
        \draw[tick] (0.013,\v) -- (-0.013,\v) node[left=1pt, rotulo] {$\rot$};
      }
      \node[below left=1pt, rotulo] at (0,0) {$0$};
      \draw[diagonal] (0,0) -- (1.05,1.05);
      \draw[curva] plot[domain=0:1, samples=100, smooth] (\x, {Gb(\x)});
      \fill[raiz] (1, 1) circle (2pt);
      \node[below right=2pt, nota] at (1, 1) {$1$};
      \node[anchor=south, font=\small, text=text] at (0.57, 1.14) {(b) $G'(1)\leqslant 1$};
    \end{scope}
  \end{tikzpicture}
  \caption{A curva $y=G(x)$ e a reta $y=x$ nos dois casos possíveis para a
    derivada em $1$. No painel (a), em que $G'(1)>1$, há duas interseções em
    $[0,1]$, e a menor delas é $e$: a reta corta a curva de baixo para cima
    antes de $x=1$, de modo que a órbita que parte de $0$ converge para $e$ e
    não para $1$. No painel (b), em que $G'(1)\leqslant 1$, a única interseção
    em $[0,1]$ é $x=1$. A reta tracejada é $y=x$; os pontos marcados são as
    raízes de $x=G(x)$.}
  \label{fig:extincao-raizes}
\end{figure}

\begin{proof}
Começamos pelo caso em que nenhuma família é vazia, porque nele a conclusão é
imediata.

\smallskip
\noindent\textbf{O caso $p_0=0$.}
Se $p_0=0$, então uma geração vazia nunca pode aparecer e, portanto, $e=0$.
Nesse caso a média é estritamente maior que $1$. De fato, a hipótese
$p_1\neq1$ garante que existe algum $k\geqslant2$ com $p_k>0$ e, como $p_0=0$,
temos $\sum_{k\geqslant2}p_k=1-p_1$; logo
\[
\mu
=\sum_{k=1}^{\infty}kp_k
\geqslant p_1+2\sum_{k=2}^{\infty}p_k
=p_1+2(1-p_1)
=2-p_1
>1 .
\]
Assim, quando $p_0=0$, vale $e\neq1$ e $\mu>1$, em acordo com o enunciado.

Suponha, de agora em diante, que $p_0>0$. Nesse caso a função auxiliar
\[
h(x)=G(x)-x,\qquad 0\leqslant x\leqslant1,
\]
é contínua em $[0,1]$, e as raízes de $x=G(x)$ em $[0,1]$ são exatamente os
zeros de $h$. Como $G(1)=1$, temos $h(1)=0$, de modo que $x=1$ é sempre raiz.
Pelo \Cref{thm:extincao-raiz}, $e$ é a menor raiz não negativa; basta então
decidir se $h$ se anula em $[0,1)$.

As propriedades de $G$ que usaremos vêm da derivação termo a termo da série de
potências no interior do seu raio de convergência, que é pelo menos $1$. Para
$0\leqslant x<1$,
\[
G'(x)=\sum_{k=1}^{\infty}kp_k\,x^{k-1}\geqslant0,
\qquad
G''(x)=\sum_{k=2}^{\infty}k(k-1)p_k\,x^{k-2}\geqslant0 .
\]
Consequentemente, $G$ é não decrescente e convexa em $[0,1]$, e $G'$ é não
decrescente em $[0,1)$. Pelo \Cref{thm:fgp-momentos}, a derivada lateral
$G'(1)=\lim_{x\uparrow1}G'(x)$ coincide com $\mu$, admitindo o valor $+\infty$.
Em qualquer caso, a monotonicidade de $G'$ fornece
\[
G'(x)\leqslant G'(1)=\mu
\quad\text{para } 0\leqslant x<1,
\]
e, portanto, $h'(x)=G'(x)-1\leqslant\mu-1$ nesse intervalo. Registramos ainda
que $h'(0)=G'(0)-1=p_1-1$, quantidade negativa, já que excluímos $p_1=1$.

\smallskip
\noindent\textbf{O caso $\mu\leqslant1$.}
Aqui $h'(x)\leqslant\mu-1\leqslant0$ para todo $0\leqslant x<1$, o que mostra
que $h$ é não crescente em $[0,1)$. Como $h$ é contínua em $1$ e $h(1)=0$,
segue que $h(x)\geqslant0$ para todo $x\in[0,1)$. De fato, se $x<1$ e $y\in(x,1)$,
a monotonicidade dá $h(x)\geqslant h(y)$ e, fazendo $y\uparrow1$, obtemos
$h(x)\geqslant h(1)=0$.

Falta excluir a possibilidade de um zero de $h$ em $[0,1)$. Seja
$x_0\in[0,1)$ com $h(x_0)=0$. Para $t\in[x_0,1)$, a monotonicidade dá
$h(t)\leqslant h(x_0)=0$ e o parágrafo anterior dá $h(t)\geqslant0$; logo $h$ é
identicamente nula em $[x_0,1)$. Em particular, $h'(t)=0$ para
$t\in(x_0,1)$, isto é, $G'(t)=1$ nesse intervalo. A derivada de $G(t)-t$ é
então nula em $(x_0,1)$, de modo que $G(t)-t$ é constante ali; como $G(1)=1$,
essa constante vale $0$ e obtemos $G(t)=t$ para todo $t\in(x_0,1)$. Temos
assim uma igualdade entre séries de potências válida em um intervalo não
degenerado e, pela unicidade dos coeficientes, $p_1=1$ e $p_k=0$ para
$k\neq1$, o que contraria a hipótese $p_1\neq1$. Concluímos que, neste caso, a
única raiz de $x=G(x)$ em $[0,1]$ é $x=1$ e, pelo \Cref{thm:extincao-raiz},
$e=1$.

\smallskip
\noindent\textbf{O caso $\mu>1$.}
Como $\lim_{x\uparrow1}G'(x)=\mu>1$, existe $a\in(0,1)$ tal que $G'(x)>1$ para
todo $x\in(a,1)$, ou seja, $h'(x)>0$ nesse intervalo. Logo $h$ é estritamente
crescente em $(a,1)$. Para $a<x<1$ e $y\in(x,1)$ temos então $h(x)<h(y)$, e o
limite de uma função crescente em $1$ é o supremo dos seus valores; como esse
limite é $h(1)=0$, concluímos que
\[
h(x)<0
\quad\text{para todo } x\in(a,1).
\]
Perto de $0$, ao contrário, $h$ decresce a partir de $h(0)=p_0>0$. Fixe
$x_1\in(a,1)$; então $h(x_1)<0<h(0)=p_0$. O teorema do valor intermediário
aplicado a $h$ no intervalo $[0,x_1]$ produz $\eta\in(0,x_1)$ com
$h(\eta)=0$, isto é, uma raiz de $x=G(x)$ no interior de $[0,1]$. Pelo
\Cref{thm:extincao-raiz}, $e\leqslant\eta<1$ e, portanto, a probabilidade de
sobrevivência $1-e$ é estritamente positiva.

\smallskip
Reunindo os dois casos e a observação inicial sobre $p_0=0$, concluímos que,
sob a hipótese $p_1\neq1$, vale $e=1$ se e somente se $\mu\leqslant1$.
\end{proof}

O critério separa, portanto, dois regimes. Quando $\mu\leqslant1$, a extinção é
certa; quando $\mu>1$, a linhagem sobrevive indefinidamente com
probabilidade $1-e>0$, e $e$ é justamente a menor raiz da equação $x=G(x)$
no intervalo $[0,1)$, cuja existência a prova acima estabelece. O critério
não calcula esse
valor: para obtê-lo é preciso resolver a equação, o que só é viável quando a
função geradora de probabilidade tem forma simples. Os dois exercícios a seguir
ilustram a situação com funções geradoras cujas raízes podem ser encontradas
explicitamente.

% Origem: G&W, exercício 9.23.
\begin{exercise}\label{exercise:extincao-geometrica}
Suponha que os tamanhos de família de um processo de ramificação tenham função
de probabilidade
\[
p_k=pq^k,\qquad k=0,1,2,\ldots,
\]
com $0<p=1-q<1$. Use o \Cref{thm:extincao-raiz} para mostrar que a
probabilidade de que o processo se extinga em algum momento é $p/q$ quando
$p\leqslant1/2$. A função geradora de probabilidade do tamanho de família é a de
\eqref{eq:fgp-geometrica-geral} e a equação $x=G(x)$ se reduz, nesse caso, a
uma equação quadrática.
\end{exercise}

% Origem: G&W, exercício 9.24.
\begin{exercise}\label{exercise:extincao-binomial}
Suponha que cada família de um processo de ramificação tenha distribuição
$\mathrm{Bin}(2,p)$, com $0\leqslant p\leqslant1$ e $q=1-p$, isto é, que cada
indivíduo deixe nenhum, um ou dois filhos com probabilidades $q^2$, $2pq$
e $p^2$,
respectivamente. Mostre que a probabilidade de extinção eventual é
\[
e=
\begin{cases}
1, & \text{se } 0\leqslant p\leqslant\dfrac{1}{2},\\[0.3cm]
\left(\dfrac{q}{p}\right)^{2}, & \text{se } \dfrac{1}{2}\leqslant p\leqslant1.
\end{cases}
\]
Compare o resultado com o \Cref{thm:extincao-criterio}, lembrando que a média
do tamanho de família é $\mu=2p$.
\end{exercise}

% Origem: G&W, seção 9.6.
\section{Problemas}\label{sec:problemas}

Os problemas a seguir ampliam o que desenvolvemos nas seções anteriores e
retomam, em vários pontos, ferramentas das notas de aula sobre esperança
condicional e sobre funções geradoras de probabilidade. Não apresentamos
soluções. Em cada enunciado indicamos, quando isso não é evidente, qual
resultado já disponível deve ser empregado.

\begin{enumerate}[label={\arabic*.}]
\item
Sejam $X_1,X_2,\ldots$ variáveis aleatórias independentes, cada uma com média
$\mu$ e variância $\sigma^2$, e seja $N$ uma variável aleatória com valores no
conjunto $\{1,2,\ldots\}$ dos inteiros positivos, independente da sequência
$(X_i)_{i\geqslant 1}$. Mostre que a soma aleatória
\[
S=X_1+X_2+\cdots+X_N
\]
tem variância dada por
\[
\operatorname{Var}(S)
=\sigma^2\,\mathbb{E}[N]+\mu^2\operatorname{Var}(N).
\]
A fórmula se estende ao caso em que $N$ assume valores em $\{0,1,2,\ldots\}$,
com a convenção de que a soma vazia é $0$; é essa extensão que usamos no
processo de ramificação.

\emph{Sugestão:} decomponha sobre os valores de $N$ e use a independência para
fatorar as esperanças; recorde a fórmula da variância de uma soma de um número
fixo de parcelas independentes.

Sejam agora $Z_0,Z_1,\ldots$ os tamanhos das gerações de um processo de
ramificação em que cada tamanho de família tem média $\mu$ e variância
$\sigma^2$. Use a identidade acima com $N=Z_{n-1}$, juntamente com a relação
\eqref{eq:recursao-geracao} e o \Cref{thm:media-geracao}, para concluir que,
para $n\geqslant 1$,
\[
\operatorname{Var}(Z_n)
=\sigma^2\mu^{n-1}+\mu^2\operatorname{Var}(Z_{n-1})
=\mu\operatorname{Var}(Z_{n-1})+\sigma^2\mu^{2n-2},
\]
e deduza, iterando essa recorrência a partir de $\operatorname{Var}(Z_0)=0$,
uma expressão para $\operatorname{Var}(Z_n)$ em termos de $\mu$, $\sigma^2$ e
$n$. Compare o resultado com a fórmula obtida no
\Cref{exercise:variancia-geracao}.

\item
Use o resultado do problema anterior para mostrar que, se
$Z_0,Z_1,\ldots$ é um processo de ramificação cujos tamanhos de família têm
média $\mu>1$ e variância $\sigma^2$, então
\[
\operatorname{Var}\!\left(\frac{Z_n}{\mu^n}\right)
\longrightarrow
\frac{\sigma^2}{\mu(\mu-1)},
\qquad n\to\infty.
\]
\emph{Sugestão:} escreva a recorrência do problema anterior para a variável
normalizada $Z_n/\mu^n$ e observe que as parcelas somadas formam uma progressão
geométrica.

\item
Usando o Teorema da Partição, apresentado nas notas de aula sobre esperança
condicional, e condicionando no valor de $Z_m$, mostre que, se
$Z_0,Z_1,\ldots$ é um processo de ramificação cuja distribuição do tamanho de
família tem média $\mu$, então
\[
\mathbb{E}[Z_mZ_n]=\mu^{n-m}\,\mathbb{E}[Z_m^2],
\qquad\text{para } m<n.
\]
\emph{Sugestão:} condicione em $\{Z_m=k\}$; as $k$ linhagens que partem dos
indivíduos da geração $m$ evoluem de modo independente e cada uma contribui
para a geração $n$ com esperança $\mu^{n-m}$, como no
\Cref{thm:media-geracao}.

\item
Seja $(Z_n:n\geqslant 0)$ um processo de ramificação com $Z_0=1$. Escreva $P_r$
para a função geradora de probabilidade do tamanho da $r$-ésima geração, de
modo que $P_r=G_r$ na notação da \Cref{sec:metodo-geradoras}. Mostre que
\[
P_n(s)=P_r\bigl(P_{n-r}(s)\bigr),
\qquad 1\leqslant r\leqslant n-1,
\]
e compare essa identidade com o \Cref{exercise:composicao-geracoes}. A leitura
probabilística é que primeiro se formam as gerações até a $r$-ésima e, a partir
de cada indivíduo aí presente, se forma um processo independente com a mesma
lei, agora com $n-r$ gerações a percorrer.

Suponha agora que o processo seja modificado de tal modo que a geração inicial
$Z_0$ tenha distribuição $\mathrm{Poisson}(\lambda)$ e que o número de filhos de
cada indivíduo tenha distribuição $\mathrm{Poisson}(\lambda)$, com todas as
escolhas independentes. Seja $H_n$ a função geradora de probabilidade do número
total de indivíduos $Z_0+Z_1+\cdots+Z_n$. Encontre uma função $f$ tal que
\[
H_n(s)=f\bigl(s\,H_{n-1}(s)\bigr).
\]
\emph{Sugestão:} o fator $s$ fora da composição registra a contribuição da
geração inicial, e a escolha poissoniana faz a função geradora de probabilidade
da geração inicial reaparecer na expressão.

\smallskip
\noindent\emph{(Oxford, 1977)}

\item
Neste problema usamos $Y_n$ para o tamanho da $n$-ésima geração e $T_n$ para o
total de indivíduos, para não sobrecarregar a notação $Z_n$ já usada nas seções
anteriores. Seja então $(Y_n:n\geqslant 0)$ um processo de ramificação com
$\mathbb{P}(Y_0=1)=1$, e seja
\[
T_n=Y_0+Y_1+\cdots+Y_{n-1}
\]
o número total de indivíduos nas $n$ primeiras gerações, contando a geração
$0$. Escreva $Q_n$ para a função geradora de probabilidade de $T_n$. Mostre
que, para $n\geqslant 2$,
\[
Q_n(s)=s\,P_1\bigl(Q_{n-1}(s)\bigr),
\]
em que $P_1$ é a função geradora de probabilidade da distribuição do tamanho de
família.

\emph{Sugestão:} o fator $s$ conta o indivíduo da geração $0$ em $T_n$.
A composição reúne os totais das $Y_1$ linhagens que partem da primeira
geração, cada uma delas com $n-1$ gerações contadas a partir de seu próprio
fundador.

\smallskip
\noindent\emph{(Oxford, 1974)}

\item
\begin{enumerate}[label=(\alph*)]
  \item Explique o que se entende pelo termo ``processo de ramificação''.
  \item Seja $Z_n$ o tamanho da $n$-ésima geração de um processo de ramificação
        em que cada tamanho de família tem função geradora de probabilidade
        $G$, e suponha $Z_0=1$. Mostre que a função geradora de probabilidade
        $G_n$ do tamanho da $n$-ésima geração satisfaz
        \[
        G_{n+1}(s)=G_n\bigl(G(s)\bigr),
        \qquad n\geqslant 1.
        \]
  \item Mostre que
        \[
        G(s)=1-\alpha(1-s)^{\beta}
        \]
        é a função geradora de probabilidade de uma variável aleatória com
        valores em $\{0,1,2,\ldots\}$
        quando $\alpha,\beta\in(0,1)$ e, para essa escolha de $G$, encontre
        $G_n$ explicitamente.
  \item Encontre a probabilidade $\mathbb{P}(Z_n=0)$, mostre que ela converge,
        quando $n\to\infty$, para $1-\alpha^{1/(1-\beta)}$, e explique por que
        isso implica que a probabilidade de extinção final do processo é
        $1-\alpha^{1/(1-\beta)}$.
\end{enumerate}

\emph{Sugestões:} em (b), use a definição de $G_n$ como iterada de $G$; em
(c), examine os coeficientes da expansão de $(1-s)^{\beta}$ e use a parte (b);
em (d), use a continuidade da probabilidade para identificar o limite de
$\mathbb{P}(Z_n=0)$ com a probabilidade de que a população se aniquile em
algum momento.

\smallskip
\noindent\emph{(Cambridge, 2001)}
\end{enumerate}

%----------------------------------------------------------------------
% BIBLIOGRAFIA
%----------------------------------------------------------------------
\bibliographystyle{plain}
\bibliography{../arquivos-aux/referencias}

\end{document}
